\begin{proof}[Proof of \cref{obj:thm:exact-response-type-game}]
Write \(\mathfrak S_n\) for the group of unit permutations.  The proof packages five finite-sum facts.

\begin{enumerate}
\item Fix a labeled procedure \(p=(\mathcal D,\widehat\tau)\).  For an assignment \(A\) and observed vector \(y\), define
\[
\mathsf{den}(A)=\sum_{\sigma\in\mathfrak S_n}\mathcal D(\sigma A),
\qquad
\mathsf{num}(A,y)=
\sum_{\sigma\in\mathfrak S_n}\mathcal D(\sigma A)\,
\widehat\tau(\sigma A,\sigma y),
\]
where \((\sigma A)_i=A_{\sigma^{-1}(i)}\) and \((\sigma y)_i=y_{\sigma^{-1}(i)}\).  Let
\[
\bar{\mathcal D}(A)=\frac{1}{|\mathfrak S_n|}
\sum_{\sigma\in\mathfrak S_n}\mathcal D(\sigma A),
\]
and let
\[
\bar{\widehat\tau}(A,y)=
\begin{cases}
\operatorname{clip}_{[-L_c/2,L_c/2]}
\!\left(\mathsf{num}(A,y)/\mathsf{den}(A)\right),
& \mathsf{den}(A)>0,\\[2mm]
0,& \mathsf{den}(A)=0.
\end{cases}
\]
Since \(L_c=\sum_a |c_a|\ge 0\), the second value lies in the estimator interval.  Reindexing the finite sum by left composition with any fixed unit permutation gives
\[
\bar{\mathcal D}(\sigma_0 A)=\bar{\mathcal D}(A),
\qquad
\bar{\widehat\tau}(\sigma_0 A,\sigma_0 y)=\bar{\widehat\tau}(A,y),
\]
so \(\bar p=(\bar{\mathcal D},\bar{\widehat\tau})\) is invariant.

For fixed \(A\) with \(\mathsf{den}(A)>0\), the weighted mean
\[
\mathsf{num}(A,y)/\mathsf{den}(A)
\]
is a convex combination of values in \([-L_c/2,L_c/2]\), hence the clipping is inactive.  Therefore, for every real target value \(\theta\),
\[
\mathsf{den}(A)\{\bar{\widehat\tau}(A,y)-\theta\}^2
\le
\sum_{\sigma\in\mathfrak S_n}
\mathcal D(\sigma A)\{\widehat\tau(\sigma A,\sigma y)-\theta\}^2
\]
by Jensen's inequality for the convex map \(t\mapsto(t-\theta)^2\), applied to the weights \(\mathcal D(\sigma A)\).  If \(\mathsf{den}(A)=0\), the left side is zero and the right side is nonnegative, so the same inequality holds.

Apply this with \(y=Y^{\mathrm{obs}}(z,A)\) and \(\theta=\tau_c(z)\), then sum over \(A\).  The identities
\[
\tau_c(\sigma z)=\tau_c(z),
\qquad
Y^{\mathrm{obs}}(\sigma z,\sigma A)=
\sigma Y^{\mathrm{obs}}(z,A)
\]
convert the right-hand side into the permutation average of labeled risks:
\[
R_n(\bar p;z)
\le
\frac{1}{|\mathfrak S_n|}
\sum_{\sigma\in\mathfrak S_n}R_n(p;\sigma z).
\]
Let \(q\) be the orbit procedure obtained by collapsing the invariant procedure \(\bar p\) to allocation and observation counts.  For this collapsed procedure, the labeled and orbit descriptions put the same total probability on each allocation-count and observation-count fiber, so summing the squared loss over those fibers gives
\[
R_n(q;\operatorname{counts}(z))=R_n(\bar p;z).
\]
Hence \(q\) is dominated by the averaged invariant representative. % lean: lossless_symmetrization

\item An invariant labeled procedure is determined exactly by allocation-count and observation-count orbits.  For an assignment \(A\in\mathcal A_K^n\) and an observed binary vector \(y\), define
\[
\operatorname{alloc}(A)_a=\#\{i:A_i=a\},
\qquad
\operatorname{succ}(A,y)_a=\#\{i:A_i=a,\ y_i=1\}.
\]
Thus \(\operatorname{alloc}(A)\in\mathcal R_{n,K}\), and, when \(\operatorname{alloc}(A)=r\), the vector \(\operatorname{succ}(A,y)\) is an element of \(\prod_{a\in\mathcal A_K}\{0,\ldots,r_a\}\).  Write
\[
\mathcal O_{\mathrm{alloc}}(r)
=
\{A'\in\mathcal A_K^n:\operatorname{alloc}(A')=r\}.
\]
For an orbit procedure \(q=(\pi,\delta)\), inflate it to a labeled procedure by
\[
\mathcal D_q(A)=
\frac{\pi_{\operatorname{alloc}(A)}}{
|\mathcal O_{\mathrm{alloc}}(\operatorname{alloc}(A))|},
\qquad
\widehat\tau_q(A,y)=
\delta(\operatorname{alloc}(A),\operatorname{succ}(A,y)).
\]
Conversely, if \(p=(\mathcal D,\widehat\tau)\) is invariant, choose one representative \(A_r^\circ\in\mathcal O_{\mathrm{alloc}}(r)\) for each allocation count \(r\), and one representative pair \((A_{r,x}^\circ,y_{r,x}^\circ)\) with \(\operatorname{alloc}(A_{r,x}^\circ)=r\) and \(\operatorname{succ}(A_{r,x}^\circ,y_{r,x}^\circ)=x\) for each observation count \(x\in\prod_{a\in\mathcal A_K}\{0,\ldots,r_a\}\).  Set
\[
\pi_p(r)=
|\mathcal O_{\mathrm{alloc}}(r)|\,\mathcal D(A_r^\circ),
\qquad
\delta_p(r,x)=\widehat\tau(A_{r,x}^\circ,y_{r,x}^\circ).
\]
These definitions are independent of the representatives: equal allocation counts are exactly equality up to a unit permutation, and equal allocation and observation counts are exactly equality up to a simultaneous unit permutation of assignment and observed outcomes.

The two constructions are inverse to each other.  Under either construction,
\[
\mathcal D(A)=
\frac{\pi_{\operatorname{alloc}(A)}}{
|\mathcal O_{\mathrm{alloc}}(\operatorname{alloc}(A))|},
\qquad
\widehat\tau(A,y)=
\delta(\operatorname{alloc}(A),\operatorname{succ}(A,y)).
\]
It remains to identify the regrouped labeled risk with the orbit risk.  For a labeled schedule \(z\), put \(m_z=\operatorname{counts}(z)\) and define the observed labeled fiber
\[
\mathcal L_z(r,x)=
\{A\in\mathcal A_K^n:
\operatorname{alloc}(A)=r,
\operatorname{succ}(A,Y^{\mathrm{obs}}(z,A))=x\}.
\]
For a table \(h\in\mathcal H(m_z,r,x)\), the number of assignments in \(\mathcal L_z(r,x)\) having \(h_{t,a}\) units of response type \(t\) assigned to arm \(a\) is
\[
\prod_{t\in\mathcal T_K}
\frac{m_{z,t}!}{\prod_{a\in\mathcal A_K}h_{t,a}!}.
\]
Summing over the feasible contingency-table fiber and dividing by the allocation-orbit cardinality
\[
|\mathcal O_{\mathrm{alloc}}(r)|=
\frac{n!}{\prod_{a\in\mathcal A_K}r_a!}
\]
gives
\[
\frac{|\mathcal L_z(r,x)|}{|\mathcal O_{\mathrm{alloc}}(r)|}
=
\frac{\prod_{a\in\mathcal A_K}r_a!}{n!}
\sum_{h\in\mathcal H(m_z,r,x)}
\prod_{t\in\mathcal T_K}
\frac{m_{z,t}!}{\prod_{a\in\mathcal A_K}h_{t,a}!}
=
P_{m_z}(x\mid r),
\]
where the last equality is the likelihood formula in \cref{obj:def:orbit-likelihood}.  The target term is also constant on the response-count orbit,
\[
\tau_c(z)=
\frac1n\sum_i\sum_{a\in\mathcal A_K}c_a t_{i,a}
=
\frac1n\sum_{t\in\mathcal T_K}m_{z,t}\sum_{a\in\mathcal A_K}c_a t_a
=\tau_c(m_z).
\]
Therefore, regrouping first by allocation counts and then by observed fibers yields
\[
\begin{aligned}
R_n(p;z)
&=
\sum_r\sum_{A\in\mathcal O_{\mathrm{alloc}}(r)}
\frac{\pi_r}{|\mathcal O_{\mathrm{alloc}}(r)|}
\{\delta(r,\operatorname{succ}(A,Y^{\mathrm{obs}}(z,A)))-\tau_c(m_z)\}^2 \\
&=
\sum_r \pi_r\sum_x
\frac{|\mathcal L_z(r,x)|}{|\mathcal O_{\mathrm{alloc}}(r)|}
\{\delta(r,x)-\tau_c(m_z)\}^2 \\
&=
\sum_r \pi_r
\sum_x P_{m_z}(x\mid r)
\{\delta(r,x)-\tau_c(m_z)\}^2
=
R_n(q;m_z).
\end{aligned}
\]
This proves the stated bijection and the exact equality of allocation-count, observation-count, and risk rules. % lean: exact_invariant_procedure_correspondence

\item For every schedule \(z\), the count vector satisfies
\[
\operatorname{counts}(z)_t=\#\{i:t_i=t\}.
\]
Thus the finite sum over labeled units may be regrouped by its response-type fibers:
\[
\sum_i\sum_a c_a t_{i,a}
=
\sum_t
\#\{i:t_i=t\}\sum_a c_a t_a.
\]
Here and in the displayed target formula, \(1/n\) denotes the real inverse \((n:\mathbb R)^{-1}\), with the convention \(0^{-1}=0\).  Scaling the regrouped identity by this factor gives
\[
\tau_c(z)
=
\frac{1}{n}\sum_i\sum_a c_a t_{i,a}
=
\frac{1}{n}\sum_t m_t\sum_a c_a t_a
=
\tau_c(m),
\qquad m=\operatorname{counts}(z).
\] % lean: tauC_eq_tauCount_scheduleCounts

\item The arm-count condition \(K\ge2\) supplies nonempty assignment and orbit spaces, while the zero clipped estimator supplies nonempty procedure classes.  Risks are nonnegative, so both minimax values in \cref{obj:def:labeled-schedule-game,obj:def:finite-orbit-game} are formed over bounded-below finite risk ranges.

First fix an orbit procedure \(q\), and inflate it to its invariant labeled representative \(p_q\).  By the risk identity above, for every labeled schedule \(z\),
\[
R_n(p_q;z)=R_n(q;\operatorname{counts}(z))
\le
\max_m R_n(q;m).
\]
Taking the worst case over \(z\) and then the infimum over \(q\) gives
\[
\rho_n(c)\le G_n(c).
\]
Conversely, fix a labeled procedure \(p\).  The lossless symmetrization step gives an orbit procedure \(q_p\) such that, for every \(z\),
\[
R_n(q_p;\operatorname{counts}(z))
\le
\frac{1}{|\mathfrak S_n|}
\sum_{\sigma\in\mathfrak S_n}R_n(p;\sigma z)
\le
\max_{z'}R_n(p;z').
\]
Every response-count vector is \(\operatorname{counts}(z)\) for some schedule \(z\), so
\[
\max_m R_n(q_p;m)\le \max_{z'}R_n(p;z').
\]
Taking the infimum over \(p\) gives \(G_n(c)\le\rho_n(c)\).  Hence
\[
\rho_n(c)=G_n(c).
\] % lean: rhoN_eq_orbitGameValue

\item Put the orbit game in finite squared-loss form with state space \(\mathcal M_{n,K}\), design space \(\mathcal R_{n,K}\), observation space
\[
\mathcal X(r)=\prod_{a\in\mathcal A_K}\{0,\ldots,r_a\},
\]
likelihood coefficient
\[
\mathsf L_{\mathrm{orb}}(m,r,x)=P_m(x\mid r),
\]
target \(\tau_c(m)\), and action interval \([-L_c/2,L_c/2]\).  The interval is nonempty because \(L_c\ge0\), and \(\mathsf L_{\mathrm{orb}}(m,r,x)=P_m(x\mid r)\ge0\) follows term by term from the factorial formula in \cref{obj:def:orbit-likelihood}.

The state set \(\mathcal M_{n,K}\) is finite, the action set is the compact interval \([-L_c/2,L_c/2]\), and squared loss is bounded and convex in the action, so the game has a saddle point \citep{wald_1950}.  Applied to these data it gives an orbit procedure \(q_{\mathrm{sad}}\) and a probability law \(\nu\) on \(\mathcal M_{n,K}\) such that
\[
R_n(q_{\mathrm{sad}};m)\le
\inf_q\max_{m'}R_n(q;m')
\quad\text{for every }m,
\]
and
\[
\inf_q\max_{m'}R_n(q;m')
\le
\mathbb E_{\nu}\!\left[R_n(q;m)\right]
\quad\text{for every orbit procedure }q.
\]
By \cref{obj:def:finite-orbit-game}, the middle value is \(G_n(c)\), giving
\[
R_n(q_{\mathrm{sad}};m)\le G_n(c),
\qquad
G_n(c)\le \mathbb E_{\nu}\!\left[R_n(q;m)\right].
\] % lean: finite_orbit_game_has_saddle
\end{enumerate}
\end{proof}

\begin{proof}[Proof of \cref{obj:thm:multiarm-strict-extension}]
\begin{enumerate}
\item For \(K=2\) and \(c=(1,-1)\), the arm-count condition required in \cref{obj:thm:exact-response-type-game} is immediate. Its value-equality conclusion, applied to this two-arm contrast, gives
\[
\rho_n(c)=G_n(c)
\]
for every \(n>0\). % lean: exact_response_type_game

\item For the same contrast,
\[
L_c=|1|+|-1|=2,
\qquad
C_0(c)=\frac{L_c^2}{4}=1 .
\]
This proves the asserted first-order constant. % lean: twoArmContrast_C0

\item Fix \(n>0\) and \(m\in\mathcal M_{n,2}\). The four two-arm response types are
\[
t_{00}=(0,0),\qquad
t_{01}=(0,1),\qquad
t_{10}=(1,0),\qquad
t_{11}=(1,1).
\]
They are pairwise distinct and every \(t\in\mathcal T_2\) is exactly one of these four types: the pair \((t_1,t_2)\) is one of
\[
(0,0),\qquad (0,1),\qquad (1,0),\qquad (1,1).
\]
Thus the orbit-target sum may be evaluated over this four-element list:
\[
\tau_c(m)
=
\frac{1}{n}\sum_{t\in\{t_{00},t_{01},t_{10},t_{11}\}}m_t\sum_{a\in\mathcal A_2}c_a t_a .
\]
For \(c=(1,-1)\), their contrast scores are
\[
\sum_a c_a t_{00,a}=0,\qquad
\sum_a c_a t_{01,a}=-1,\qquad
\sum_a c_a t_{10,a}=1,\qquad
\sum_a c_a t_{11,a}=0 .
\]
Therefore
\[
\tau_c(m)
=
\frac{m_{t_{10}}-m_{t_{01}}}{n}
=
\frac{p_+(m)-p_-(m)}{n},
\]
which is the claimed two-arm target formula. % lean: twoArm_tauCount_eq_effect_difference

\item The four functions above differ by evaluating them at one of the two arms: \(t_{00}\) differs from \(t_{01}\) at the second arm, from \(t_{10}\) at the first arm, and from \(t_{11}\) at either arm; \(t_{01}\) and \(t_{10}\) differ at the first arm; \(t_{01}\) and \(t_{11}\) differ at the first arm; and \(t_{10}\) and \(t_{11}\) differ at the second arm. Hence they are pairwise distinct. % lean: twoArm_responseTypes_pairwise_distinct

\item If \(t\in\mathcal T_2\), then the pair \((t_1,t_2)\) is one of
\[
(0,0),\qquad (0,1),\qquad (1,0),\qquad (1,1),
\]
so \(t\) is respectively \(t_{00},t_{01},t_{10}\), or \(t_{11}\). Thus these four types exhaust \(\mathcal T_2\). % lean: twoArm_responseType_cases

\item With these four response types kept as the elements of \(\mathcal T_2\), the likelihood identity is exactly the specialization of \cref{obj:def:orbit-likelihood} to \(K=2\):
\[
P_m(x\mid r)
=
\frac{\prod_{a\in\mathcal A_2} r_a!}{n!}
\sum_{h\in\mathcal H(m,r,x)}
\prod_{t\in\mathcal T_2}
\frac{m_t!}{\prod_{a\in\mathcal A_2} h_{t,a}!}.
\]
Here the fiber \(\mathcal H(m,r,x)\) imposes the response-type margins \(m\), allocation margins \(r\), and arm-success totals \(x\), as in \cref{obj:def:orbit-likelihood}. % lean: orbitLik

\item Let \(n\ge 8\). Applying \cref{obj:thm:embedded-two-arm-converse} with \(K=2\) and \(c=(1,-1)\) gives
\[
d_n(c)\le 43\,C_0(c)n^{-4/3}.
\]
Using \(C_0(c)=1\) from the calculation above yields
\[
d_n(c)\le 43\,n^{-4/3}.
\]
% lean: embedded_two_arm_converse

\item For the lower second-order bound, apply \cref{obj:thm:universal-second-order-rate} with \(K=2\) and \(c=(1,-1)\). Let \(\kappa=\kappa_c>0\), and choose \(N_\star\) large enough that the shrinkage procedure in that result satisfies, for every \(n\ge N_\star\),
\[
\sup_{z\in\mathcal T_2^n}
R_n(\mathcal D^\star,\widehat\tau_n^{\mathrm{sh}};z)
\le
C_0(c)\{n^{-1}-\kappa n^{-4/3}\}.
\]
Since \(\rho_n(c)\) is the infimum of the same worst-case risk over all admissible procedures,
\[
\rho_n(c)
\le
C_0(c)\{n^{-1}-\kappa n^{-4/3}\}.
\]
With \(C_0(c)=1\) and \(d_n(c)=n^{-1}-\rho_n(c)\), this rearranges to
\[
\kappa n^{-4/3}\le d_n(c)
\]
for every \(n\ge N_\star\). % lean: twoArm_secondOrder_lower_witness

\item Now fix \(K\ge3\), \(n\), and a contrast \(c\). Since \(K\ge3\) implies \(K\ge2\), \cref{obj:thm:exact-response-type-game} applies and gives
\[
\rho_n(c)=G_n(c).
\]
The response-type set is the set of binary functions on \(\mathcal A_K\). Because \(|\mathcal A_K|=K\), its cardinality is
\[
|\mathcal T_K|=|\{0,1\}^{\mathcal A_K}|=2^K .
\]
% lean: exact_response_type_game

\item Fix \(N,L,U\), a finite design \(\mathcal D\) on two-arm assignments, and a measurable Hull estimator \(\widehat\tau\). Set
\[
v=\sup_Y R_{\mathrm{Hull}}(\mathcal D,\widehat\tau;Y),
\]
where \(Y\) ranges over Hull schedules with parameters \(N,L,U\). Choosing
\[
\mathcal D'=\mathcal D,
\qquad
\widehat\tau'=\widehat\tau
\]
preserves measurability and gives
\[
v=\sup_Y R_{\mathrm{Hull}}(\mathcal D',\widehat\tau';Y).
\]
Thus \(v\) belongs to the displayed risk-value class. % lean: multiarm_strict_extension

\item Finally fix \(K\ge3\), \(n>0\), and \(L<U\). First choose a number
\[
v_{\mathrm{sep}}\in [L,U],
\qquad
v_{\mathrm{sep}}\ne 0,
\qquad
v_{\mathrm{sep}}\ne 1 .
\]
The choice is by endpoint cases. If \(L=0\) and \(U=1\), take \(v_{\mathrm{sep}}=1/2\). If \(L=0\) and \(U\ne1\), then the subcase \(U=0\) contradicts \(L<U\), and in the remaining subcase take \(v_{\mathrm{sep}}=U\). If \(L\ne0\) and \(L=1\), then the subcases \(U=0\) and \(U=1\) contradict \(L<U\), and in the remaining subcase take \(v_{\mathrm{sep}}=U\). If \(L\ne0\) and \(L\ne1\), take \(v_{\mathrm{sep}}=L\). In every surviving case this gives the displayed membership and exclusions.

Define the Hull schedule
\[
Y_i(a)=v_{\mathrm{sep}},
\qquad
i=1,\ldots,n,\quad a\in\mathcal A_2 .
\]
This schedule lies in \(\mathcal Y^{\mathrm{Hull}}_{n,L,U}\). Since \(n>0\), choose one unit \(i_0\). If this embedded Hull schedule belonged to the binary \(K\)-arm schedule domain, then at the first arm its value would be the value of a binary potential outcome, hence either \(0\) or \(1\). But the same coordinate equals \(v_{\mathrm{sep}}\), contradicting \(v_{\mathrm{sep}}\notin\{0,1\}\). Therefore
\[
\mathcal Y^{\mathrm{Hull}}_{n,L,U}\nsubseteq \mathcal T_K^n .
\]
% lean: hull_not_subset_binary
\end{enumerate}
\end{proof}

\begin{proof}[Proof of \cref{obj:thm:first-order-saddle}]
\begin{enumerate}
\item We first prove the limiting assertion. The hypotheses on \(c\) imply \(K\ge 2\): if \(K=0\), there are no arm coefficients and the contrast is the zero function; if \(K=1\), the zero-sum identity forces the unique coefficient to be zero. Both alternatives contradict the assumed nonzero contrast. % lean: contrast_has_admissible_arm_count

For \(n\ge 1\), set
\[
d_n(c)=\frac{C_0(c)}{n}-\rho_n(c).
\]
With \(K\ge 2\), \(n\ge 1\), and \(c\) nonzero and zero-sum, \cref{obj:thm:embedded-two-arm-converse} applies. In particular, for every \(n\ge 8\),
\[
0\le d_n(c),
\qquad
d_n(c)\le 43\,C_0(c)n^{-4/3}.
\]
Since \(C_0(c)=L_c^2/4\ge 0\), for \(n\ge 8\) we have
\[
n\rho_n(c)-C_0(c)
=
n\left\{\frac{C_0(c)}{n}-d_n(c)\right\}-C_0(c)
=
-nd_n(c),
\]
and therefore
\[
\left|n\rho_n(c)-C_0(c)\right|
=
nd_n(c)
\le
43\,C_0(c)n^{-1/3}.
\]
The upper bound tends to \(0\), so the squeeze theorem gives
\[
\lim_{n\to\infty} n\rho_n(c)=C_0(c).
\]
% lean: firstOrder_minimax_limit

\item We next prove the finite-sample risk bound. Fix \(n\ge 1\) and a labeled schedule \(z\in\mathcal T_K^n\). Since \(c\) is nonzero,
\[
L_c=\sum_{a\in\mathcal A_K}|c_a|>0.
\]
Thus \(q_a^\star=|c_a|/L_c\) is a probability vector, and \(q_a^\star>0\) exactly on the active support \(S_c=\{a:c_a\ne0\}\).

For each unit \(i\) and arm \(a\), define the centered one-unit score on all of \(\mathcal A_K\) by
\[
\psi_i(a;z)=
\begin{cases}
\dfrac{c_a(t_{i,a}-1/2)}{q_a^\star}, & q_a^\star>0,\\[1.2ex]
0, & q_a^\star=0.
\end{cases}
\]
Under the independent \(q^\star\) design in \cref{obj:def:first-order-procedure}, the unprojected estimator is
\[
\widetilde\tau_{\mathrm{raw}}(A,z)
=
\frac1n\sum_{i=1}^n \psi_i(A_i;z).
\]
For each \(i\),
\[
\mathbb E\{\psi_i(A_i;z)\}
=
\sum_{a\in S_c} q_a^\star
\frac{c_a(t_{i,a}-1/2)}{q_a^\star}
=
\sum_{a\in\mathcal A_K} c_a(t_{i,a}-1/2)
=
\sum_{a\in\mathcal A_K} c_a t_{i,a},
\]
where the last equality uses \(\sum_a c_a=0\). Hence
\[
\mathbb E\{\widetilde\tau_{\mathrm{raw}}(A,z)\}
=
\frac1n\sum_{i=1}^n\sum_{a\in\mathcal A_K}c_at_{i,a}
=
\tau_c(z).
\]
% lean: firstOrderUnitScore_mean

Also, since \(t_{i,a}\in\{0,1\}\),
\[
(t_{i,a}-1/2)^2=\frac14,
\]
and therefore
\[
\mathbb E\{\psi_i(A_i;z)^2\}
=
\frac14\sum_{a\in S_c}\frac{c_a^2}{q_a^\star}
=
\frac14\sum_{a\in S_c}L_c|c_a|
=
\frac{L_c^2}{4}
=
C_0(c).
\]
Thus
\[
\operatorname{Var}\{\psi_i(A_i;z)\}
=
C_0(c)-\{\mathbb E\psi_i(A_i;z)\}^2
\le C_0(c).
\]
Independence across units gives
\[
\operatorname{MSE}\{\widetilde\tau_{\mathrm{raw}};z\}
=
\operatorname{Var}\{\widetilde\tau_{\mathrm{raw}}(A,z)\}
=
\frac1{n^2}\sum_{i=1}^n
\operatorname{Var}\{\psi_i(A_i;z)\}
\le
\frac{C_0(c)}{n}.
\]
% lean: firstOrderUnitScore_secondMoment

It remains to pass from the raw rule to the projected rule. For any two zero-sum contrasts \(\mathfrak c_1,\mathfrak c_2\), put
\[
\operatorname{dist}(\mathfrak c_1,\mathfrak c_2)
=
\frac12\sum_{a\in\mathcal A_K}
|(\mathfrak c_1)_a-(\mathfrak c_2)_a|.
\]
For any binary response type, splitting the arms according to \(t_a=1\) and \(t_a=0\), the zero-sum identity for \(\mathfrak c_1-\mathfrak c_2\) yields
\[
\left|
\sum_{a\in\mathcal A_K}
\{(\mathfrak c_1)_a-(\mathfrak c_2)_a\}t_a
\right|
\le
\operatorname{dist}(\mathfrak c_1,\mathfrak c_2).
\]
Averaging over the \(n\) units gives
\[
|\tau_{\mathfrak c_1}(z)-\tau_{\mathfrak c_2}(z)|
\le
\operatorname{dist}(\mathfrak c_1,\mathfrak c_2).
\]
Taking \(\mathfrak c_1=c\) and \(\mathfrak c_2=-c\), we have
\[
\tau_{-c}(z)=-\tau_c(z),
\qquad
\operatorname{dist}(c,-c)=L_c,
\]
so
\[
|\tau_c(z)|\le \frac{L_c}{2},
\qquad\text{that is,}\qquad
\tau_c(z)\in[-L_c/2,L_c/2].
\]
% lean: tauC_mem_naturalInterval

For every \(y\in[-L_c/2,L_c/2]\) and every real \(x\),
\[
\left\{
\operatorname{proj}_{[-L_c/2,L_c/2]}(x)-y
\right\}^2
\le
(x-y)^2.
\]
Indeed, if \(x\) lies in the interval there is equality; if \(x<-L_c/2\), then
\(0\le y+L_c/2\le y-x\); and if \(x>L_c/2\), then
\(0\le L_c/2-y\le x-y\). Squaring gives the displayed inequality in the two boundary cases. Applying it with \(y=\tau_c(z)\) shows, pointwise in the assignment,
\[
\left\{
\widehat\tau_{\mathrm{cHT}}^\star(A,Y^{\mathrm{obs}})
-\tau_c(z)
\right\}^2
\le
\left\{
\widetilde\tau_{\mathrm{raw}}(A,z)-\tau_c(z)
\right\}^2.
\]
Consequently,
\[
R_n(\mathcal D^\star,\widehat\tau_{\mathrm{cHT}}^\star;z)
\le
\frac{C_0(c)}{n}.
\]
Since \(z\) was arbitrary,
\[
\max_{z\in\mathcal T_K^n}
R_n(\mathcal D^\star,\widehat\tau_{\mathrm{cHT}}^\star;z)
\le
\frac{C_0(c)}{n}.
\]
% lean: contrastWeightedProcedure_risk

\item The two preceding steps are exactly the two components of the asserted conjunction: the first gives the asymptotic value of the unrestricted labeled-schedule minimax risk, and the second gives the finite-sample worst-case bound for the independent contrast-weighted design together with the projected contrast-weighted Horvitz--Thompson rule. % lean: first_order_saddle
\end{enumerate}
\end{proof}

\begin{proof}[Proof of \cref{obj:thm:embedded-two-arm-converse}]
Write
\[
R_n(\mathcal D,\widehat\tau;z)
=
\mathbb E_{\mathcal D}\!\left[
\{\widehat\tau((A_i,Y_i^{\mathrm{obs}})_{i=1}^n)-\tau_c(z)\}^2
\right]
\]
for the schedulewise risk appearing in \cref{obj:def:labeled-schedule-game}. For the two-arm contrast \((1,-1)\), write
\[
\tau_2(y)=\frac1n\sum_{i=1}^n(y_{i,+}-y_{i,-}),
\qquad
R_{n,2}(\mathcal D_2,\widehat\upsilon;y)
=
\mathbb E_{\mathcal D_2}\!\left[
\{\widehat\upsilon((B_i,V_i)_{i=1}^n)-\tau_2(y)\}^2
\right],
\]
for \(y\in\mathcal T_2^n\), with \(y_{i,+},y_{i,-}\in\{0,1\}\).

\begin{enumerate}
\item Define the sign groups and scale
\[
S_c^+=\{a\in\mathcal A_K:c_a>0\},\qquad
S_c^-=\{a\in\mathcal A_K:c_a<0\},\qquad
S_c^0=\{a\in\mathcal A_K:c_a=0\},
\qquad
h_c=\frac{L_c}{2}.
\]
Since \(c\) is nonzero and zero-sum, both \(S_c^+\) and \(S_c^-\) are nonempty, and
\[
\sum_{a\in S_c^+}c_a
=
-\sum_{a\in S_c^-}c_a
=
\frac12\sum_{a\in\mathcal A_K}|c_a|
=
h_c,
\qquad
h_c^2=C_0(c).
\]
For \(y\in\mathcal T_2^n\), embed it into a \(K\)-arm schedule \(E_cy\in\mathcal T_K^n\) by
\[
(E_cy)_{i,a}
=
\begin{cases}
y_{i,+},&a\in S_c^+,\\
y_{i,-},&a\in S_c^-,\\
0,&a\in S_c^0.
\end{cases}
\]
Then
\[
\tau_c(E_cy)
=
\frac1n\sum_{i=1}^n
\left\{
y_{i,+}\sum_{a\in S_c^+}c_a+
y_{i,-}\sum_{a\in S_c^-}c_a
\right\}
=
h_c\,\tau_2(y).
\]

Fix an arbitrary \(K\)-arm procedure \((\mathcal D,\widehat\tau)\). Let
\(\Omega=\{0,1\}^{\{1,\ldots,n\}}\), put the product law \(\mathcal L\) on
\((A,\omega)\in\mathcal A_K^n\times\Omega\) by drawing \(A\sim\mathcal D\) and independent fair bits
\(\omega_i\), and define the induced two-arm assignment
\[
\sigma_{c,i}(A,\omega)
=
\begin{cases}
+,&A_i\in S_c^+,\\
-,&A_i\in S_c^-,\\
+,&A_i\in S_c^0\text{ and }\omega_i=1,\\
-,&A_i\in S_c^0\text{ and }\omega_i=0.
\end{cases}
\]
For \(v\in\{0,1\}^n\), define the inactive-zero observation vector
\[
\widetilde v_i(A,v)
=
\begin{cases}
v_i,&c_{A_i}\ne0,\\
0,&c_{A_i}=0.
\end{cases}
\]
Let \(\mathcal D_2^\circ\) be the pushforward of \(\mathcal L\) under \(\sigma_c\), so that
\[
\mathcal D_2^\circ\{b\}
=
\mathcal L\{(A,\omega):\sigma_c(A,\omega)=b\},
\qquad b\in\{+,-\}^n.
\]
All conditional means below use the following totalization convention: on a cell with
\(\mathcal D_2^\circ\{b\}=0\), the conditional mean is set to \(0\); this value is immaterial in risks because the cell has zero \(\mathcal D_2^\circ\)-mass. Define the induced two-arm estimator, on every \(b\in\{+,-\}^n\) and \(v\in\{0,1\}^n\), by
\[
\widehat\upsilon^\circ(b,v)
=
\begin{cases}
\mathbb E_{\mathcal L}\!\left[
h_c^{-1}\widehat\tau(A,\widetilde v(A,v))
\,\middle|\, \sigma_c(A,\omega)=b
\right],&\mathcal D_2^\circ\{b\}>0,\\[1ex]
0,&\mathcal D_2^\circ\{b\}=0.
\end{cases}
\]
Because \(\widehat\tau\) is clipped to \([-h_c,h_c]\), this conditional mean lies in \([-1,1]\). For each \(y\in\mathcal T_2^n\), expanding the risk under the pushed-forward law and applying conditional Jensen's inequality gives
\[
\begin{aligned}
R_{n,2}(\mathcal D_2^\circ,\widehat\upsilon^\circ;y)
&=
\mathbb E_{\mathcal D_2^\circ}\!\left[
\{\widehat\upsilon^\circ(B,(y_{i,B_i})_{i=1}^n)-\tau_2(y)\}^2
\right]\\
&\le
\mathbb E_{\mathcal L}\!\left[
\{h_c^{-1}\widehat\tau(A,\widetilde v(A,(y_{i,\sigma_{c,i}(A,\omega)})_{i=1}^n))-\tau_2(y)\}^2
\right] \\
&=
h_c^{-2}
R_n(\mathcal D,\widehat\tau;E_cy).
\end{aligned}
\]
The last equality follows from the observation identity
\[
\widetilde v_i(A,(y_{j,\sigma_{c,j}(A,\omega)})_{j=1}^n)=(E_cy)_{i,A_i},
\qquad i=1,\ldots,n,
\]
and the target identity \(\tau_c(E_cy)=h_c\tau_2(y)\). Thus the construction supplies the statewise comparison
\[
R_{n,2}(\mathcal D_2^\circ,\widehat\upsilon^\circ;y)
\le
h_c^{-2}R_n(\mathcal D,\widehat\tau;E_cy),
\qquad y\in\mathcal T_2^n.
\]
Taking the supremum over \(y\), and using \(E_cy\in\mathcal T_K^n\), yields the induced-procedure worst-case bound
\[
\sup_{y\in\mathcal T_2^n}R_{n,2}(\mathcal D_2^\circ,\widehat\upsilon^\circ;y)
\le
h_c^{-2}\sup_{z\in\mathcal T_K^n}R_n(\mathcal D,\widehat\tau;z).
\]
Since \(\rho_{n,2}\) is the infimum of the left-hand worst-case risk over all two-arm procedures,
\[
\rho_{n,2}
\le
h_c^{-2}\sup_{z\in\mathcal T_K^n}R_n(\mathcal D,\widehat\tau;z).
\]
Multiplying by \(h_c^2=C_0(c)\) and taking the infimum over the arbitrary \(K\)-arm procedure yields
\[
C_0(c)\rho_{n,2}\le \rho_n(c).
\]
% lean: CausalSmith.Experimentation.MultiarmSecondorderMinimaxFrontier.embeddedTwoArmLowerBound

\item The identity \(C_0(c)=L_c^2/4\) gives \(C_0(c)\ge0\). Consider the product design and projected estimator in \cref{obj:def:first-order-procedure}. For every schedule \(z\in\mathcal T_K^n\), introduce the auxiliary totalized weights
\[
\bar q_a^\star=\frac{|c_a|}{L_c},
\qquad a\in\mathcal A_K.
\]
These agree with the probabilities \(q_a^\star\) from \cref{obj:def:first-order-procedure} on \(S_c\), and \(\bar q_a^\star=0\) exactly when \(c_a=0\). For \(a\in\mathcal A_K\), define
\[
\psi_i(a;z)
=
\begin{cases}
\dfrac{c_a\{z_{i,a}-1/2\}}{\bar q_a^\star},&\bar q_a^\star>0,\\[1.2ex]
0,&\bar q_a^\star=0.
\end{cases}
\]
Under the independent assignment law in \cref{obj:def:first-order-procedure}, equivalently with totalized weights \(\bar q^\star\) on \(\mathcal A_K\),
\[
\mathbb E\{\psi_i(A_i;z)\}
=
\sum_{a\in S_c} \bar q_a^\star
\frac{c_a(z_{i,a}-1/2)}{\bar q_a^\star}
=
\sum_a c_a(z_{i,a}-1/2)
=
\sum_a c_az_{i,a},
\]
because \(\sum_a c_a=0\). Also,
\[
\mathbb E\{\psi_i(A_i;z)^2\}
=
\sum_{a\in S_c}
\bar q_a^\star
\frac{c_a^2(z_{i,a}-1/2)^2}{(\bar q_a^\star)^2}
=
\frac14\sum_{a\in S_c}\frac{c_a^2}{\bar q_a^\star}
=
\frac14\sum_{a\in S_c}L_c|c_a|
=
C_0(c).
\]
Thus, for
\[
\overline\psi(A;z)=\frac1n\sum_{i=1}^n\psi_i(A_i;z),
\]
the product structure gives
\[
\mathbb E\{\overline\psi(A;z)\}=\tau_c(z),
\qquad
\operatorname{Var}\{\overline\psi(A;z)\}
=
\frac1{n^2}\sum_{i=1}^n\operatorname{Var}\{\psi_i(A_i;z)\}
\le
\frac{C_0(c)}{n}.
\]
For each unit,
\[
-\frac{L_c}{2}\le \sum_a c_a z_{i,a}\le \frac{L_c}{2},
\]
so \(\tau_c(z)\in[-L_c/2,L_c/2]\). Projection onto this interval cannot increase squared distance from \(\tau_c(z)\), hence
\[
R_n(\mathcal D^\star,\widehat\tau_{\mathrm{cHT}}^\star;z)
\le
\operatorname{Var}\{\overline\psi(A;z)\}
\le
\frac{C_0(c)}{n}.
\]
Taking the maximum over \(z\) and then the infimum over procedures gives
\[
\rho_n(c)\le \frac{C_0(c)}{n}.
\]
% lean: CausalSmith.Experimentation.MultiarmSecondorderMinimaxFrontier.contrastWeightedProcedure_upperRisk

\item Suppose \(|S_c|=2\). Let \(a_c^+\) and \(a_c^-\) be the unique active arms with positive and negative coefficients. The preceding sign-sum identities give
\[
c_{a_c^+}=h_c,\qquad c_{a_c^-}=-h_c,
\qquad c_a=0\quad(a\notin\{a_c^+,a_c^-\}).
\]
For any two-arm procedure \((\mathcal D_2,\widehat\upsilon)\), lift a two-arm assignment \(b\in\{+,-\}^n\) to the \(K\)-arm assignment \(A^b\) given by
\[
A^b_i=
\begin{cases}
a_c^+,&b_i=+,\\
a_c^-,&b_i=-.
\end{cases}
\]
Let \(\mathcal D_2^\uparrow\) be the pushforward of \(\mathcal D_2\) under \(b\mapsto A^b\). For a general \(K\)-arm assignment \(A\), define the associated two-arm assignment
\[
b^A_i=
\begin{cases}
+,&A_i=a_c^+,\\
-,&A_i\ne a_c^+,
\end{cases}
\]
and define the lifted estimator on its full domain by
\[
\widehat\tau^\uparrow(A,v)=h_c\widehat\upsilon(b^A,v),
\qquad A\in\mathcal A_K^n,
\quad v\in\{0,1\}^n.
\]
For a \(K\)-arm schedule \(z\), define the active two-arm schedule
\[
z^\downarrow_{i,+}=z_{i,a_c^+},
\qquad
z^\downarrow_{i,-}=z_{i,a_c^-}.
\]
The coefficient identities above give
\[
\tau_c(z)=h_c\tau_2(z^\downarrow).
\]
Moreover \(b^{A^b}=b\) and the lifted observations satisfy
\[
(z^\downarrow)_{i,b_i}=z_{i,A^b_i},
\qquad i=1,\ldots,n.
\]
Therefore, for every \(z\in\mathcal T_K^n\), the lifted procedure has the statewise risk identity
\[
R_n(\mathcal D_2^\uparrow,\widehat\tau^\uparrow;z)
=
h_c^2 R_{n,2}(\mathcal D_2,\widehat\upsilon;z^\downarrow).
\]
It follows that
\[
\sup_{z\in\mathcal T_K^n}R_n(\mathcal D_2^\uparrow,\widehat\tau^\uparrow;z)
\le
h_c^2\sup_{y\in\mathcal T_2^n}R_{n,2}(\mathcal D_2,\widehat\upsilon;y).
\]
Taking the infimum over two-arm procedures and using \(h_c^2>0\) to pull the fixed scale through the infimum gives the minimax comparison
\[
\rho_n(c)
\le
h_c^2\rho_{n,2}
=
C_0(c)\rho_{n,2}.
\]
Together with the lower bound obtained above,
\[
\rho_n(c)=C_0(c)\rho_{n,2}
\qquad\text{whenever } |S_c|=2.
\]
% lean: CausalSmith.Experimentation.MultiarmSecondorderMinimaxFrontier.supportTwoUpperBound

\item Now assume \(n\ge8\). From the upper bound,
\[
d_n(c)=\frac{C_0(c)}{n}-\rho_n(c)\ge0.
\]
From the lower bound,
\[
d_n(c)
=
C_0(c)n^{-1}-\rho_n(c)
\le
C_0(c)n^{-1}-C_0(c)\rho_{n,2}
=
C_0(c)\{n^{-1}-\rho_{n,2}\}.
\]
Finally, \cref{obj:thm:coarse-two-arm-minimax-lower} gives
\[
n^{-1}-43n^{-4/3}\le \rho_{n,2},
\]
and therefore
\[
n^{-1}-\rho_{n,2}\le 43n^{-4/3}.
\]
Since \(C_0(c)\ge0\),
\[
d_n(c)
\le
C_0(c)\{n^{-1}-\rho_{n,2}\}
\le
43\,C_0(c)n^{-4/3}.
\]
% lean: CausalSmith.Experimentation.MultiarmSecondorderMinimaxFrontier.embedded_two_arm_converse

\item In the support-two case, the equality established above gives
\[
d_n(c)
=
\frac{C_0(c)}{n}-\rho_n(c)
=
C_0(c)n^{-1}-C_0(c)\rho_{n,2}
=
C_0(c)\{n^{-1}-\rho_{n,2}\}.
\]
This proves both asserted identities for \(|S_c|=2\).
% lean: CausalSmith.Experimentation.MultiarmSecondorderMinimaxFrontier.embedded_two_arm_converse
\end{enumerate}
\end{proof}

\begin{proof}[Proof of \cref{obj:thm:universal-second-order-rate}]
Fix \(K\ge 2\) and a nonzero zero-sum contrast \(c\). We prove the five asserted conclusions in the order displayed.

\begin{enumerate}
\item Let
\[
\sigma_c(t)=\frac{\sum_{a\in\mathcal A_K}c_a t_a}{L_c/2},
\qquad t\in\mathcal T_K .
\]
Since \(c\neq 0\), choose \(a_0\) with \(c_{a_0}\neq 0\). For the response type \(t^{a_0}\) with \(t^{a_0}_{a_0}=1\) and \(t^{a_0}_a=0\) for \(a\neq a_0\), the numerator is \(c_{a_0}\), so the finite set of nonzero values \(|\sigma_c(t)|\) is nonempty and has a positive minimum. This is exactly \(\lambda_c\), hence \(0<\lambda_c\). Also, zero-summability gives
\[
c_{a_0}=-\sum_{a\neq a_0}c_a,
\qquad
|c_{a_0}|\le \sum_{a\neq a_0}|c_a|,
\]
and therefore \(2|c_{a_0}|\le L_c\). Thus one admissible normalized nonzero value is at most one, and the minimum satisfies \(\lambda_c\le 1\). Finally,
\[
A_c=\frac{\sqrt{\lambda_c}}2-\frac{\lambda_c}{16}>0,
\qquad
B_c=\frac{7\lambda_c}{16}>0,
\]
because \(0<\lambda_c\le 1\). Hence
\[
\kappa_c=\frac12\min\{A_c,B_c\}>0 .
\]
% lean: lambdaC_pos_le_one, kappaC_pos

\item The first-order rule \(\widehat\tau_{\mathrm{cHT}}^\star\) from \cref{obj:def:first-order-procedure} is controlled first before projection. For the degenerate case \(n=0\), there is a single empty schedule, its target is
\[
\tau_c(\varnothing)=0,
\]
and the centered contrast-weighted sum is the empty sum, hence equals \(0\). The projection fixes \(0\), so the first-order procedure has risk
\[
R_0(\mathcal D^\star,\widehat\tau_{\mathrm{cHT}}^\star;\varnothing)=0.
\]
Since squared-error risks are nonnegative, this gives \(\rho_0(c)=0\). The reciprocal estimate now begins with positive population size, so assume \(n\ge 1\). For \(z=(t_i)_{i=1}^n\), define, for an assignment vector \(A\in\mathcal A_K^n\),
\[
\xi_i(A_i,z)
=
\sum_{a\in S_c}
 c_a\mathbf 1\{A_i=a\}\frac{t_{i,a}-1/2}{q_a^\star},
\qquad
\widetilde\tau_n(A,z)=\frac1n\sum_{i=1}^n\xi_i(A_i,z).
\]
Under the independent \(q^\star\)-design, inverse-probability cancellation and \(\sum_a c_a=0\) give
\[
\mathbb E\{\xi_i(A_i,z)\mid z\}
=
\sum_{a\in\mathcal A_K} c_a t_{i,a},
\qquad
\mathbb E\{\xi_i(A_i,z)^2\mid z\}=C_0(c),
\]
and hence
\[
\mathbb E\{\widetilde\tau_n(A,z)\mid z\}=\tau_c(z),
\qquad
\operatorname{Var}\{\widetilde\tau_n(A,z)\mid z\}\le \frac{C_0(c)}{n}.
\]
The target lies in the projection interval: for each response type \(t\), the positive and negative coefficient masses are both \(L_c/2\), so
\[
-\frac{L_c}{2}
\le
\sum_{a\in\mathcal A_K}c_a t_a
\le
\frac{L_c}{2},
\qquad
-\frac{L_c}{2}\le \tau_c(z)\le \frac{L_c}{2}.
\]
For every \(y\in[-L_c/2,L_c/2]\), interval projection satisfies
\[
\left|\operatorname{proj}_{[-L_c/2,L_c/2]}(x)-y\right|\le |x-y|,
\]
because the projection fixes \(x\) inside the interval and otherwise moves \(x\) to the nearer endpoint between \(x\) and \(y\). Therefore
\[
\left\{\operatorname{proj}_{[-L_c/2,L_c/2]}(\widetilde\tau_n(A,z))-\tau_c(z)\right\}^2
\le
\left\{\widetilde\tau_n(A,z)-\tau_c(z)\right\}^2.
\]
Taking expectations gives
\[
R_n(\mathcal D^\star,\widehat\tau_{\mathrm{cHT}}^\star;z)\le \frac{C_0(c)}{n}
\qquad\text{for every }z\in\mathcal T_K^n,
\]
and taking the infimum over procedures gives \(\rho_n(c)\le C_0(c)/n\) for \(n\ge1\). Together with the zero-unit branch, this proves the first-order envelope for every \(n\).
% lean: minimaxEnvelopeBound_of_contrastWeightedProcedure

\item Since
\[
j_{\mathrm{tail}}\exp\{-j_{\mathrm{tail}}^{1/3}/8\}\longrightarrow 0
\qquad(j_{\mathrm{tail}}\to\infty),
\]
multiplying by \(4\sqrt{\lambda_c}\) gives
\[
4\sqrt{\lambda_c}\,j_{\mathrm{tail}}\exp\{-j_{\mathrm{tail}}^{1/3}/8\}\longrightarrow 0 .
\]
Because \(A_c>0\), there is \(N_0\) such that, for all integers \(j_{\mathrm{tail}}\ge N_0\),
\[
4\sqrt{\lambda_c}\,j_{\mathrm{tail}}\exp\{-j_{\mathrm{tail}}^{1/3}/8\}\le A_c .
\]
Set \(N_c=\max\{N_0,1\}\). Then the stated tail side condition holds for every integer threshold variable at least \(N_c\).
% lean: shrinkage_tail_prefactor_tendsto_zero

For the risk bound, fix \(n\ge N_c\) and a schedule \(z=(t_i)_{i=1}^n\). Under the independent \(q^\star\)-design, define
\[
V_i=\frac{1}{L_c/2}
\sum_{a\in S_c}
c_a\mathbf 1\{A_i=a\}\frac{t_{i,a}-1/2}{q_a^\star},
\qquad
X_n=\frac1n\sum_{i=1}^n V_i,
\]
\[
\Theta_n(z)=\frac1n\sum_{i=1}^n\sigma_c(t_i),
\qquad
Q_n(z)=\sum_{i=1}^n\sigma_c(t_i)^2 .
\]
Then
\[
\mathbb E(X_n\mid z)=\Theta_n(z),
\qquad
\operatorname{Var}(X_n\mid z)=\frac1n-\frac{Q_n(z)}{n^2},
\]
and the definition of \(\lambda_c\) gives
\[
Q_n(z)\ge \lambda_c\sum_{i=1}^n|\sigma_c(t_i)|
\ge \lambda_c n|\Theta_n(z)|.
\]
For the tail estimate, write
\[
\zeta_i(A)=V_i(A)-\sigma_c(t_i),
\qquad
X_n(A)-\Theta_n(z)=\frac1n\sum_{i=1}^n\zeta_i(A).
\]
The product design makes the variables \(\zeta_i\) independent conditional on \(z\). Also \(V_i(A)\in[-1,1]\) and \(\mathbb E(V_i\mid z)=\sigma_c(t_i)\), so \(\zeta_i\) has conditional mean zero and lies in the interval \([-1-\sigma_c(t_i),1-\sigma_c(t_i)]\), whose length is \(2\). Hoeffding's lemma therefore gives, for every real \(\alpha_{\mathrm{mgf}}\),
\[
\mathbb E\!\left[\exp\{\alpha_{\mathrm{mgf}}\zeta_i\}\mid z\right]
\le
\exp\{\alpha_{\mathrm{mgf}}^2/2\}.
\]
Multiplying the moment-generating-function bounds across the independent units yields
\[
\mathbb E\!\left[\exp\!\left\{\alpha_{\mathrm{mgf}}\sum_{i=1}^n\zeta_i\right\}\mid z\right]
\le
\exp\{n\alpha_{\mathrm{mgf}}^2/2\}.
\]
Thus Chernoff's bound, with \(\alpha_{\mathrm{mgf}}=\omega_{\mathrm{tail}}\) for \(\omega_{\mathrm{tail}}>0\) and with the endpoint \(\omega_{\mathrm{tail}}=0\) immediate, gives
\[
\Pr\!\left\{\sum_{i=1}^n\zeta_i\ge n\omega_{\mathrm{tail}}\mid z\right\}
\le \exp\{-n\omega_{\mathrm{tail}}^2/2\}
\qquad(\omega_{\mathrm{tail}}\ge0).
\]
Applying the same argument to \(-\sum_i\zeta_i\) and taking the union of the two one-sided events gives, for every \(\omega_{\mathrm{tail}}\ge0\),
\[
\Pr\{|X_n-\Theta_n(z)|\ge \omega_{\mathrm{tail}}\mid z\}
\le 2\exp\{-n\omega_{\mathrm{tail}}^2/2\}.
\]
% lean: normalizedScore_mean_variance, normalizedTypeScore_sq_lower, normalizedScore_tail_bound

Put
\[
b_n=n^{-1/3},
\qquad
\varepsilon_n=\frac{\sqrt{\lambda_c}}4 b_n,
\qquad
g_n(x)=\max\{-b_n,\min\{x,b_n\}\}.
\]
The clipped-shrinkage estimator is
\[
\widehat\tau_n^{\mathrm{sh}}
=
\frac{L_c}{2}\{X_n-\varepsilon_n g_n(X_n)\},
\]
followed by projection onto \([-L_c/2,L_c/2]\); that projection fixes the displayed value. Indeed, fix a realized assignment and write \(x=X_n\). The bounds
\[
|x|\le 1,\qquad 0\le b_n\le 1,\qquad 0\le\varepsilon_n\le \frac14\le 1
\]
imply
\[
|x-\varepsilon_n g_n(x)|\le 1.
\]
To see this, use the three regions in the definition of \(g_n\). If \(x<-b_n\), then \(g_n(x)=-b_n\) and
\[
-1\le x\le x+\varepsilon_n b_n\le x+b_n<0\le 1.
\]
If \(b_n<x\), then \(g_n(x)=b_n\) and
\[
-1\le 0<x-b_n\le x-\varepsilon_n b_n\le x\le 1.
\]
In the remaining case, \(-b_n\le x\le b_n\), so \(g_n(x)=x\) and
\[
|x-\varepsilon_n g_n(x)|=(1-\varepsilon_n)|x|\le |x|\le 1.
\]
Thus, for every realized assignment,
\[
-\frac{L_c}{2}
\le
\frac{L_c}{2}\{X_n-\varepsilon_n g_n(X_n)\}
\le
\frac{L_c}{2},
\]
because \(L_c/2\ge0\), and the final interval projection is inactive.
% lean: shrunkenClippedScore_abs_le_one
Let
\[
\mathcal C_n(z)=\operatorname{Cov}\{X_n,g_n(X_n)\mid z\}.
\]
Since \(g_n\) is monotone,
\[
\mathcal C_n(z)\ge 0,
\]
and
\[
\mathbb E\!\left[\{X_n-\varepsilon_n g_n(X_n)-\Theta_n(z)\}^2\mid z\right]
\le
\operatorname{Var}(X_n\mid z)-2\varepsilon_n\mathcal C_n(z)+\varepsilon_n^2b_n^2 .
\]
% lean: finiteDesign_cov_monotone_comp_nonneg, finiteDesign_mse_shrinkage_upper

If \(|\Theta_n(z)|\le b_n/2\), define the far-tail event
\[
\mathscr E_{\mathrm{tail}}(z)
=
\left\{A\in\mathcal A_K^n:
|X_n(A)-\Theta_n(z)|\ge b_n/2\right\}.
\]
The tail bound gives
\[
\Pr\{\mathscr E_{\mathrm{tail}}(z)\mid z\}
\le 2\exp\{-n^{1/3}/8\}.
\]
On \(\mathscr E_{\mathrm{tail}}(z)^c\), the local condition gives
\[
|X_n(A)|
\le |X_n(A)-\Theta_n(z)|+|\Theta_n(z)|
\le b_n,
\]
so \(g_n(X_n(A))=X_n(A)\). Since \(|X_n(A)|\le1\) and \(b_n\le1\), clipping also gives the uniform perturbation bound
\[
|g_n(X_n(A))-X_n(A)|\le 1.
\]
Combining the two displays,
\[
|g_n(X_n(A))-X_n(A)|
\le \mathbf 1\{A\in\mathscr E_{\mathrm{tail}}(z)\}.
\]
Moreover \(|X_n(A)-\Theta_n(z)|\le |X_n(A)|+|\Theta_n(z)|\le2\). Since \(\mathbb E(X_n\mid z)=\Theta_n(z)\),
\[
\mathcal C_n(z)
=
\operatorname{Var}(X_n\mid z)
+\mathbb E\!\left[(X_n-\Theta_n(z))\{g_n(X_n)-X_n\}\mid z\right],
\]
and therefore
\[
\begin{aligned}
\mathcal C_n(z)
&\ge \operatorname{Var}(X_n\mid z)
-\mathbb E\!\left[|X_n-\Theta_n(z)|\,|g_n(X_n)-X_n|\mid z\right]\\
&\ge \operatorname{Var}(X_n\mid z)-2\Pr\{\mathscr E_{\mathrm{tail}}(z)\mid z\}\\
&\ge \operatorname{Var}(X_n\mid z)-4\exp\{-n^{1/3}/8\}.
\end{aligned}
\]
Since \(0\le\varepsilon_n\le1/4\),
\[
1-2\varepsilon_n\ge0,
\qquad
Q_n(z)\ge0.
\]
Substituting \(\operatorname{Var}(X_n\mid z)=n^{-1}-Q_n(z)/n^2\) and the covariance lower bound into the preceding MSE inequality gives
\[
\begin{aligned}
&\mathbb E\!\left[\{X_n-\varepsilon_n g_n(X_n)-\Theta_n(z)\}^2\mid z\right] \\
&\qquad\le
\left(1-2\varepsilon_n\right)
\left\{\frac1n-\frac{Q_n(z)}{n^2}\right\}
+8\varepsilon_n\exp\{-n^{1/3}/8\}
+\varepsilon_n^2b_n^2 \\
&\qquad\le
\frac1n-\frac{2\varepsilon_n}{n}
+8\varepsilon_n\exp\{-n^{1/3}/8\}
+\varepsilon_n^2b_n^2,
\end{aligned}
\]
where the last line drops the nonpositive term
\[-(1-2\varepsilon_n)Q_n(z)/n^2\]. Using
\[
\frac{b_n}{n}=n^{-4/3},
\qquad
\varepsilon_n=\frac{\sqrt{\lambda_c}}4 b_n,
\qquad
\varepsilon_n^2b_n^2
=
\frac{\lambda_c}{16}n^{-4/3},
\]
we obtain
\[
\mathbb E\!\left[\{X_n-\varepsilon_n g_n(X_n)-\Theta_n(z)\}^2\mid z\right]
\le
\frac1n-A_c n^{-4/3}
+2\sqrt{\lambda_c}\,b_n\exp\{-n^{1/3}/8\}.
\]
The choice of \(N_c\) implies
\[
2\sqrt{\lambda_c}\,b_n\exp\{-n^{1/3}/8\}
\le
\frac{A_c}{2}n^{-4/3},
\]
so the local case gives
\[
\mathbb E\!\left[\{X_n-\varepsilon_n g_n(X_n)-\Theta_n(z)\}^2\mid z\right]
\le
\frac1n-\frac{A_c}{2}n^{-4/3}
\le
\frac1n-\kappa_c n^{-4/3}.
\]
% lean: shrinkageProcedure_risk_bound_of_tail

If \(|\Theta_n(z)|>b_n/2\), then the lower bound on \(Q_n(z)\) yields
\[
\frac{Q_n(z)}{n^2}\ge \frac{\lambda_c}{2}n^{-4/3}.
\]
Using \(\mathcal C_n(z)\ge 0\) and the same identity for \(\varepsilon_n^2b_n^2\),
\[
\mathbb E\!\left[\{X_n-\varepsilon_n g_n(X_n)-\Theta_n(z)\}^2\mid z\right]
\le
\frac1n-\frac{7\lambda_c}{16}n^{-4/3}
\le
\frac1n-\kappa_c n^{-4/3}.
\]
The two cases cover every schedule. Multiplying by \((L_c/2)^2=C_0(c)\) gives, uniformly in \(z\),
\[
R_n(\mathcal D^\star,\widehat\tau_n^{\mathrm{sh}};z)
\le
C_0(c)\{n^{-1}-\kappa_c n^{-4/3}\}.
\]
Taking the supremum over schedules proves the asserted clipped-shrinkage risk bound for all \(n\ge N_c\).
% lean: shrinkageProcedure_risk_bound_of_tail

\item Since \(\rho_n(c)\) is the minimax infimum, the shrinkage bound implies, for all \(n\ge N_c\),
\[
\rho_n(c)
\le
C_0(c)\{n^{-1}-\kappa_c n^{-4/3}\}.
\]
Thus
\[
d_n(c)=\frac{C_0(c)}n-\rho_n(c)
\ge
C_0(c)\kappa_c n^{-4/3},
\]
and therefore
\[
n^{4/3}d_n(c)\ge C_0(c)\kappa_c
\qquad(n\ge N_c).
\]
On the other hand, \cref{obj:thm:embedded-two-arm-converse} gives, for all \(n\ge 8\),
\[
d_n(c)\le 43C_0(c)n^{-4/3},
\]
so
\[
n^{4/3}d_n(c)\le 43C_0(c)
\qquad(n\ge 8).
\]
The eventual lower and upper bounds imply
\[
C_0(c)\kappa_c
\le
\liminf_{n\to\infty} n^{4/3}d_n(c)
\le
\limsup_{n\to\infty} n^{4/3}d_n(c)
\le
43C_0(c).
\]
% lean: universal_second_order_rate, embedded_two_arm_converse, secondOrderScale_mul_inverse

\item Let \(a_n\) be a positive regularly varying sequence with index \(\beta\), and suppose
\[
a_n d_n(c)\longrightarrow C
\qquad\text{with }C>0.
\]
Define the auxiliary positive sequence
\[
\upsilon_n=
\begin{cases}
1,& n=0,\\
a_n n^{-4/3},& n\ge 1.
\end{cases}
\]
For \(n\ge1\),
\[
\upsilon_n\, n^{4/3}d_n(c)=a_nd_n(c).
\]
The bounds from the previous step and the convergence to \(C\) imply that, eventually,
\[
\frac{C}{2\cdot 43C_0(c)}
\le
\upsilon_n
\le
\frac{2C}{C_0(c)\kappa_c}.
\]
Regular variation at the scale factor \(2\) gives
\[
\frac{a_{2n}}{a_n}\longrightarrow 2^\beta,
\]
and hence
\[
\frac{\upsilon_{2n}}{\upsilon_n}
=
\frac{a_{2n}}{a_n}\,2^{-4/3}
\longrightarrow
2^{\beta-4/3}.
\]
A positive sequence that is eventually bounded above and away from zero can have a dyadic ratio limit only equal to \(1\): indeed, for the dyadic logarithms
\[
\chi_k=\log \upsilon_{2^k},
\]
the increments \(\chi_{k+1}-\chi_k\) converge to the logarithm of the dyadic ratio, while the eventual boundedness of \(\upsilon_{2^k}\) makes \(\chi_k\) eventually bounded; Cesaro averaging of
\[
\frac{\chi_J-\chi_0}{J}
=
\frac1J\sum_{k=0}^{J-1}(\chi_{k+1}-\chi_k)
\]
therefore forces that logarithm to be \(0\). Applying this to \(\upsilon_n\) gives
\[
2^{\beta-4/3}=1.
\]
Since the map \(r\mapsto 2^r\) is strictly increasing,
\[
\beta=\frac43 .
\]
% lean: dyadic_ratio_limit_eq_one, universal_second_order_rate
\end{enumerate}
\end{proof}

\begin{proof}[Proof of \cref{obj:thm:coarse-two-arm-minimax-lower}]
Fix \(n\ge 8\) and put
\[
\alpha:=n^{-1/3}.
\]
\begin{enumerate}
\item Since \(n>0\), also \(\alpha>0\), and the real-power identities give
\[
\alpha^3=n^{-1},\qquad
n^{-2/3}=\alpha^2,\qquad
n^{2/3}=n\alpha .
\]
Moreover \(\alpha\le 1/2\). Indeed, \(n^{-1}\le 8^{-1}\). If \(1/2<\alpha\), then
\[
\frac14<\alpha^2
\]
and, multiplying by the positive inequality \(1/2<\alpha\),
\[
\frac18<\alpha^3=n^{-1},
\]
contradicting \(n^{-1}\le 8^{-1}\). The same identities also give
\[
\frac1{\alpha^2}=n^{2/3}.
\]
% lean: coarse_two_arm_minimax_lower

\item We first prove the sharper fractional lower bound in the auxiliary parameter \(\alpha\). Let
\[
w_\alpha(\theta)
=
\frac{15}{8\alpha}
\left(1-\frac{4\theta^2}{\alpha^2}\right)^2
\mathbf 1\{|\theta|\le \alpha/2\}.
\]
This is a probability density on \((-1/2,1/2)\), is positive on its support interior, and vanishes at the support endpoints. Direct differentiation on \(|\theta|<\alpha/2\) gives
\[
w_\alpha'(\theta)
=
-\frac{30\theta}{\alpha^3}
\left(1-\frac{4\theta^2}{\alpha^2}\right),
\]
so
\[
J(w_\alpha)
=
\int \frac{w_\alpha'(\theta)^2}{w_\alpha(\theta)}\,d\theta
=
\int_{-\alpha/2}^{\alpha/2}\frac{480\theta^2}{\alpha^5}\,d\theta
=
\frac{40}{\alpha^2}.
\]
For fixed \(\theta\) with \(|\theta|\le \alpha/2\), draw independent two-arm response types \(r_i=(r_{i,1},r_{i,2})\in\{0,1\}^2\) with
\[
\Pr_\theta\{(1,0)\}=\frac{\alpha+\theta}{2},\qquad
\Pr_\theta\{(0,1)\}=\frac{\alpha-\theta}{2},\qquad
\Pr_\theta\{(0,0)\}=\Pr_\theta\{(1,1)\}=\frac{1-\alpha}{2}.
\]
These masses are nonnegative and sum to one because \(0<\alpha\le 1/2\) and \(|\theta|\le \alpha/2\). Let
\[
\sigma_i:=r_{i,1},\qquad
X(\sigma):=\sum_{i=1}^n \mathbf 1\{\sigma_i=1\},\qquad
\bar R(\sigma):=\frac1n\sum_{i=1}^n(2\mathbf 1\{\sigma_i=1\}-1).
\]
Then the marginal law of \(\sigma\) is the product Bernoulli law
\[
p_\theta(\sigma)
=
\prod_{i=1}^n
\left(\frac{1+\theta}{2}\right)^{\mathbf 1\{\sigma_i=1\}}
\left(\frac{1-\theta}{2}\right)^{\mathbf 1\{\sigma_i=0\}},
\]
with Fisher information
\[
I_n(\theta)=\frac{n}{1-\theta^2}
\le
\frac{n}{1-\alpha^2/4}.
\]
The finite-population contrast target for a response vector \(r\) is
\[
\tau(r)=\frac1n\sum_{i=1}^n(r_{i,1}-r_{i,2}).
\]
Conditional on \(\sigma_i=1\),
\[
\mathbb E_\theta(r_{i,1}-r_{i,2}\mid \sigma_i=1)=\frac{\alpha+\theta}{1+\theta},
\]
and conditional on \(\sigma_i=0\),
\[
\mathbb E_\theta(r_{i,1}-r_{i,2}\mid \sigma_i=0)=-\frac{\alpha-\theta}{1-\theta}.
\]
Thus the posterior target is
\[
\gamma_\alpha(\theta,\sigma)
:=
\mathbb E_\theta(\tau(r)\mid \sigma)
=
\theta+\omega_\alpha(\theta)\{\bar R(\sigma)-\theta\},
\qquad
\omega_\alpha(\theta):=\frac{\alpha-\theta^2}{1-\theta^2}.
\]
For \(|\theta|\le\alpha/2\),
\[
0\le \omega_\alpha(\theta)\le \alpha.
\]
At fixed \(\sigma\),
\[
\partial_\theta\gamma_\alpha(\theta,\sigma)
=
1-\omega_\alpha(\theta)
+
\frac{2\theta(\alpha-1)}{(1-\theta^2)^2}
\{\bar R(\sigma)-\theta\}.
\]
Since \(\mathbb E_\theta\bar R(\sigma)=\theta\),
\[
\mathbb E_\theta\{\partial_\theta\gamma_\alpha(\theta,\sigma)\}
=
1-\omega_\alpha(\theta)
\ge
1-\alpha.
\]
For any scalar estimator \(t:\{0,\ldots,n\}\to\mathbb R\), let \(t^\circ\) be its clipping to \([-1,1]\). Since \(\tau(r)\in[-1,1]\),
\[
\{t^\circ(X(\sigma))-\tau(r)\}^2
\le
\{t(X(\sigma))-\tau(r)\}^2.
\]
The conditional square completion gives, for each \(\theta\) and \(\sigma\),
\[
p_\theta(\sigma)\{t^\circ(X(\sigma))-\gamma_\alpha(\theta,\sigma)\}^2
\le
\sum_r
\Pr_\theta(r)\mathbf 1\{r_{i,1}=\sigma_i\ \forall i\}
\{t^\circ(X(\sigma))-\tau(r)\}^2 .
\]
Applying \cref{obj:lem:bayesian-information-inequality} to the finite observation space \(\{0,1\}^n\), the density \(w_\alpha\), likelihood \(p_\theta\), target \(\gamma_\alpha\), and estimator \(t^\circ(X(\sigma))\), with the regularity hypotheses discharged by the displayed compact support and finite product likelihood, yields
\[
\int w_\alpha(\theta)
\sum_\sigma
p_\theta(\sigma)
\{t^\circ(X(\sigma))-\gamma_\alpha(\theta,\sigma)\}^2\,d\theta
\ge
\frac{(1-\alpha)^2}
{40/\alpha^2+n/(1-\alpha^2/4)}.
\]
Combining this estimate with the square-completion and clipping inequalities shows that every scalar estimator has Bayes risk at least
\[
\frac{(1-\alpha)^2}
{40/\alpha^2+n/(1-\alpha^2/4)}.
\]
Let \(\nu_\alpha\) be the induced distribution on triples
\[
\vartheta=(p_+,p_-,r_0),
\]
where \(p_+\) counts \((1,0)\) units, \(p_-\) counts \((0,1)\) units, and \(r_0\) counts equal-potential-outcome units. Applied to \(\nu_\alpha\), \cref{obj:lem:two-arm-scalar-prior-schedule-kernel} gives a complete-schedule prior whose Rao--Blackwell scalar risk is dominated by every design-estimator Bayes risk. Hence, by \cref{obj:def:labeled-schedule-game},
\[
\frac{(1-\alpha)^2}
{40/\alpha^2+n/(1-\alpha^2/4)}
\le
\rho_{n,2}.
\]
Using \(n^{-2/3}=\alpha^2\), \(\alpha=n^{-1/3}\), and \(1/\alpha^2=n^{2/3}\), this is exactly
\[
\frac{(1-n^{-1/3})^2}
{n/(1-n^{-2/3}/4)+40n^{2/3}}
\le
\rho_{n,2}.
\]
% lean: twoArmPosteriorCompat_fraction_le_rho2

\item It remains to convert the fractional bound to the advertised coarse form. Since \(0<\alpha\le 1/2\),
\[
0<1-\frac{\alpha^2}{4}
\]
and
\[
\left(1-\frac{\alpha^2}{4}\right)^{-1}\le 1+\alpha.
\]
Therefore
\[
40\alpha+\left(1-\frac{\alpha^2}{4}\right)^{-1}
\le
1+41\alpha.
\]
We claim
\[
1-43\alpha
\le
\frac{(1-\alpha)^2}
{40\alpha+(1-\alpha^2/4)^{-1}}.
\]
The denominator is positive. If \(1-43\alpha\ge 0\), then
\[
(1-43\alpha)\left\{40\alpha+\left(1-\frac{\alpha^2}{4}\right)^{-1}\right\}
\le
(1-43\alpha)(1+41\alpha)
\le
(1-\alpha)^2.
\]
If \(1-43\alpha<0\), then the left-hand side after multiplication by the positive denominator is at most \(0\), while \((1-\alpha)^2\ge0\). This proves the claim in both cases.

Finally,
\[
n^{-4/3}=n^{-1}\alpha
\]
and
\[
\frac{n}{1-\alpha^2/4}+40n\alpha
=
n\left\{40\alpha+\left(1-\frac{\alpha^2}{4}\right)^{-1}\right\}.
\]
Thus
\[
\begin{aligned}
n^{-1}-43n^{-4/3}
&=
n^{-1}(1-43\alpha)\\
&\le
n^{-1}
\frac{(1-\alpha)^2}
{40\alpha+(1-\alpha^2/4)^{-1}}\\
&=
\frac{(1-\alpha)^2}
{n/(1-\alpha^2/4)+40n\alpha}.
\end{aligned}
\]
The last expression is the fractional lower bound proved in the preceding step, so
\[
n^{-1}-43n^{-4/3}\le \rho_{n,2}.
\]
Together with the first displayed bound, this proves both conclusions.
% lean: coarseTwoArmNumericalBound
\end{enumerate}
\end{proof}

\begin{proof}[Proof of \cref{obj:prop:k3-lp-certificate}]
Let \(n,M\ge 1\).  We write \(\rho_n^\dagger=\rho_n(c^\dagger)\) and use the identification of \(\Lambda_{n,M}^\dagger\) with the three-arm specialization of the rational grid value from \cref{obj:def:k3-rational-lp}.

\begin{enumerate}
\item Applying \cref{obj:thm:rational-contrast-grid-certificate-sandwich} with \(K=3\), contrast \(c^\dagger\), and mesh \(M\) gives certificates \(\pi^{n,M,c^\dagger}\), \(w^{n,M,c^\dagger}\), \(u\), \(\nu_{n,M,c^\dagger}\), and a barycenter estimator \(\delta_{n,M,c^\dagger}\) satisfying
\[
B_{n,c^\dagger}(\nu_{n,M,c^\dagger})
\le
\rho_n^\dagger
\le
U_{n,M}(c^\dagger)
\le
\Lambda_{n,M}^\dagger
\le
B_{n,c^\dagger}(\nu_{n,M,c^\dagger})
+\frac{C_0(c^\dagger)}{4M^2}.
\]
For \(c^\dagger=(1,-1/2,-1/2)\),
\[
L_{c^\dagger}=|1|+|-1/2|+|-1/2|=2,
\qquad
C_0(c^\dagger)=\frac{L_{c^\dagger}^2}{4}=1 .
\]
The displayed sandwich first gives
\[
\rho_n^\dagger\le \Lambda_{n,M}^\dagger .
\]
It also gives
\[
\Lambda_{n,M}^\dagger
\le
B_{n,c^\dagger}(\nu_{n,M,c^\dagger})+\frac{1}{4M^2}
\le
\rho_n^\dagger+\frac{1}{4M^2}.
\]
Since \(M\ge1\),
\[
\frac{1}{4M^2}\le \frac{2}{M}+\frac{1}{M^2},
\]
and hence
\[
\Lambda_{n,M}^\dagger
\le
\rho_n^\dagger+\frac{2}{M}+\frac{1}{M^2}.
\]
% lean: CausalSmith.Experimentation.MultiarmSecondorderMinimaxFrontier.k3_lp_certificate

\item The unrestricted minimax risk is an infimum of suprema of squared-error risks, so
\[
0\le \rho_n^\dagger .
\]
Combining this with \(\rho_n^\dagger\le\Lambda_{n,M}^\dagger\) yields
\[
0\le\Lambda_{n,M}^\dagger .
\]
% lean: CausalSmith.Experimentation.MultiarmSecondorderMinimaxFrontier.k3_lp_certificate

\item It remains to record the elementary feasible point giving the upper bound \(1\).  Let \(r^\star=(n,0,0)\in\mathcal R_{n,3}\), let \(0\in\Gamma_M\) be the grid action corresponding to the index \(j=M\), and define
\[
\pi_r=\mathbf 1\{r=r^\star\},
\qquad
w_{r,x,g}=\mathbf 1\{r=r^\star,\ g=0\},
\qquad
u_\star=1 .
\]
The normalization constraints in \cref{obj:def:k3-rational-lp} are then immediate.  For a response-count vector \(m\in\mathcal M_{n,3}\), set
\[
\alpha(\zeta)=\zeta_1-\frac12\zeta_2-\frac12\zeta_3,
\qquad \zeta\in\mathcal T_3 .
\]
Because each \(\zeta_a\in\{0,1\}\), one has \(-1\le\alpha(\zeta)\le1\).  Therefore
\[
\tau_{c^\dagger}(m)
=
\frac1n\sum_{\zeta\in\mathcal T_3}m_\zeta\,\alpha(\zeta)
\in[-1,1],
\qquad
\tau_{c^\dagger}(m)^2\le1 .
\]
Using the orbit likelihood of \cref{obj:def:orbit-likelihood}, whose finite probabilities in \(x\) are nonnegative and sum to \(1\), the risk row for the displayed feasible point is
\[
\sum_{r,x,g}P_m(x\mid r)w_{r,x,g}
\{g-\tau_{c^\dagger}(m)\}^2
=
\sum_x P_m(x\mid r^\star)\tau_{c^\dagger}(m)^2
=
\tau_{c^\dagger}(m)^2
\le 1 .
\]
Thus \(u_\star=1\) is feasible, and the minimizing value satisfies
\[
\Lambda_{n,M}^\dagger\le1 .
\]
% lean: CausalSmith.Experimentation.MultiarmSecondorderMinimaxFrontier.gridLPValueK3_le_one

\item Apply \cref{obj:thm:rational-contrast-grid-certificate-sandwich} again, now with mesh \(M=n^2\).  Since \(n\ge1\), also \(n^2\ge1\).  The exact rational certificate clause gives a rational allocation mass \(\pi\), real joint design-action mass \(w\), rational epigraph value \(u\), and rational response-count multipliers \(\nu\), packaged as an exact primal-dual certificate.  Unpacking that certificate supplies rational grid weights \(w^{\mathbb Q}\) and dual multipliers \(y\) such that
\[
w=w^{\mathbb Q}\ \text{after real embedding},\qquad
(\pi,w^{\mathbb Q},u)\ \text{is primal feasible},\qquad
y\ \text{is dual feasible},
\]
\[
\nu=y_{\mathrm{risk}},
\qquad
u=\operatorname{obj}_{\mathrm{dual}}(y).
\]
The optimal-value clause gives
\[
u=\Lambda_{n,n^2}^\dagger .
\]
% lean: CausalSmith.Experimentation.MultiarmSecondorderMinimaxFrontier.k3_lp_certificate

\item Fix \(t\in\mathcal T_3\).  Choose any labeling of the \(n\) units by \(\{1,\ldots,n\}\), and let \(z^{(t)}\in\mathcal T_3^n\) be the constant labeled schedule
\[
z_i^{(t)}=t,\qquad i=1,\ldots,n,
\]
and let \(m^{(t)}\in\mathcal M_{n,3}\) be its response-count vector,
\[
m^{(t)}_\zeta=\#\{i:z_i^{(t)}=\zeta\},
\qquad \zeta\in\mathcal T_3 .
\]
Then \(m^{(t)}_t=n\).  Since \(n\ge1\),
\[
m^{(t)}_t>0,
\]
so the response type \(t\) is represented by a response-count vector with total population size \(n\).
% lean: CausalSmith.Experimentation.MultiarmSecondorderMinimaxFrontier.k3_lp_certificate

\item Finally, work at the second-order scale \(k^{4/3}\).  The asserted limit follows from the following direct calculation.  For all \(k\ge1\),
\[
k^{4/3}\left(\frac{2}{k^2}+\frac{1}{k^4}\right)
=
2k^{-2/3}+k^{-8/3}.
\]
Both powers on the right tend to zero as \(k\to\infty\), because \(2/3>0\) and \(8/3>0\).  Hence
\[
\lim_{k\to\infty}
k^{4/3}\left(\frac{2}{k^2}+\frac{1}{k^4}\right)=0 .
\]
% lean: CausalSmith.Experimentation.MultiarmSecondorderMinimaxFrontier.k3_lp_certificate
\end{enumerate}
\end{proof}

\begin{proof}[Proof of \cref{obj:thm:rational-contrast-grid-certificate-sandwich}]
\begin{enumerate}
\item Since a zero-sum contrast on \(K=0\) or \(K=1\) is identically zero, the hypotheses imply \(K\geq2\).  The rational grid program of \cref{obj:def:rational-contrast-grid-lp} is a finite rational linear program.  A feasible point is obtained by placing all allocation mass on one allocation orbit and all grid-action mass at \(0\), with an epigraph coordinate large enough to dominate the finitely many risk rows.  The risk rows also force every feasible epigraph coordinate to be nonnegative.  Finite-dimensional rational LP strong duality therefore gives rational primal and dual optimizers.  Decoding the primal optimizer gives
\[
\pi^{n,M,c},\qquad w^{\mathbb Q},\qquad u\in\mathbb Q,
\]
and decoding the dual optimizer gives multipliers \(y\), whose risk-row component is
\[
\nu_{n,M,c}(m)=y_m^{\mathrm{risk}},
\qquad m\in\mathcal M_{n,K}.
\]
The decoded certificate satisfies
\[
w^{n,M,c}_{r,x,g}=w^{\mathbb Q}_{r,x,g},
\qquad
u=y^{\mathrm{norm},-}-y^{\mathrm{norm},+},
\qquad
u=\Lambda_{n,M}(c),
\]
with primal feasibility and dual feasibility in the grid program.  In particular,
\[
y_m^{\mathrm{risk}}\geq0
\quad\text{for all }m,
\qquad
\sum_{m\in\mathcal M_{n,K}}y_m^{\mathrm{risk}}=1,
\]
so \(\nu_{n,M,c}\) is a rational prior on response-count vectors. % lean: exists_exact_grid_primal_dual_certificate

\item Define the estimator, for \(r\in\mathcal R_{n,K}\) and compatible \(x\), by
\[
\delta_{n,M,c}(r,x)=
\begin{cases}
\displaystyle
\sum_{g\in\Gamma_{M,c}} g\,w_{r,x,g}^{n,M,c}/\pi_r^{n,M,c},
& \pi_r^{n,M,c}>0,\\[0.7em]
0,& \pi_r^{n,M,c}=0.
\end{cases}
\]
This is exactly the barycenter condition in the certificate statement; the zero convention applies precisely on allocation orbits with zero certificate mass. % lean: IsGridBarycenter

\item For a rational prior \(\nu\) on \(\mathcal M_{n,K}\), set
\[
D_{r,x,c}(\nu)
=
\sum_{m\in\mathcal M_{n,K}}\nu_m P_m(x\mid r),
\qquad
N_{r,x,c}(\nu)
=
\sum_{m\in\mathcal M_{n,K}}\nu_m P_m(x\mid r)\tau_c(m),
\]
and use the posterior-mean convention
\[
\bar\tau_{r,x,c}(\nu)=
\begin{cases}
N_{r,x,c}(\nu)/D_{r,x,c}(\nu),&D_{r,x,c}(\nu)>0,\\
0,&D_{r,x,c}(\nu)=0.
\end{cases}
\]
Then the Bayes risk at allocation orbit \(r\) is
\[
\operatorname{BayesRisk}_{n,c}(r,\nu)
=
\sum_m\nu_m\tau_c(m)^2
-
\sum_x
\begin{cases}
N_{r,x,c}(\nu)^2/D_{r,x,c}(\nu),&D_{r,x,c}(\nu)>0,\\
0,&D_{r,x,c}(\nu)=0.
\end{cases}
\]
Completing the square gives, for every invariant orbit procedure \((\pi,\delta)\),
\[
\sum_m\nu_m R_n((\pi,\delta);m)
\geq
\sum_r\pi_r\,\operatorname{BayesRisk}_{n,c}(r,\nu)
\geq
B_{n,c}(\nu).
\]
Thus \(B_{n,c}(\nu)\) is a Bayes lower bound for the finite orbit game.  By the value equality in \cref{obj:thm:exact-response-type-game},
\[
B_{n,c}(\nu)\leq \rho_n(c).
\]
Applying this to \(\nu=\nu_{n,M,c}\) gives the stated lower certificate inequality. % lean: lowerCertificate_le_rhoN

\item The barycenter rule \(\delta_{n,M,c}\), together with \(\pi^{n,M,c}\), is an admissible orbit procedure.  Through the invariant correspondence and value equality in \cref{obj:thm:exact-response-type-game},
\[
\rho_n(c)
\leq
\sup_{m\in\mathcal M_{n,K}}
\sum_r \pi_r^{n,M,c}
\sum_x P_m(x\mid r)
\{\delta_{n,M,c}(r,x)-\tau_c(m)\}^2
=
U_{n,M}(c).
\]
This proves \(\rho_n(c)\leq U_{n,M}(c)\). % lean: rhoN_le_upperCertificate

\item For \(\pi_r^{n,M,c}>0\), let
\[
\kappa_{r,x}^{n,M,c}(g)
=
w_{r,x,g}^{n,M,c}/\pi_r^{n,M,c},
\qquad g\in\Gamma_{M,c}.
\]
Primal feasibility gives \(\sum_g\kappa_{r,x}^{n,M,c}(g)=1\), and the barycenter definition gives
\[
\delta_{n,M,c}(r,x)
=
\sum_g g\,\kappa_{r,x}^{n,M,c}(g).
\]
Hence, for every \(m,r,x\) with \(\pi_r^{n,M,c}>0\),
\[
\{\delta_{n,M,c}(r,x)-\tau_c(m)\}^2
\leq
\sum_g\kappa_{r,x}^{n,M,c}(g)\{g-\tau_c(m)\}^2.
\]
When \(\pi_r^{n,M,c}=0\), the corresponding contribution is multiplied by zero.  Multiplying by \(\pi_r^{n,M,c}P_m(x\mid r)\), summing over \((r,x)\), and using the primal risk row gives
\[
U_{n,M}(c)\leq u=\Lambda_{n,M}(c).
\]
 % lean: upperCertificate_le_gridObjective

\item It remains to compare \(\Lambda_{n,M}(c)\) with the lower certificate.  Fix an allocation orbit \(r\).  Since each \(\tau_c(m)\) lies in \([-h_c,h_c]\), the posterior mean \(\bar\tau_{r,x,c}(\nu_{n,M,c})\) also lies in this interval when \(D_{r,x,c}(\nu_{n,M,c})>0\), and the convention value \(0\) lies there when \(D_{r,x,c}(\nu_{n,M,c})=0\).  The grid \(\Gamma_{M,c}\) has mesh \(h_c/M\), so choose \(\gamma^\star_{r,x}\in\Gamma_{M,c}\) with
\[
\left|\gamma^\star_{r,x}-\bar\tau_{r,x,c}(\nu_{n,M,c})\right|
\leq \frac{h_c}{2M}.
\]
Write the relevant dual occupancy difference and nonnegativity multipliers as
\[
\alpha_{r,x}=y^{\mathrm{occ},+}_{r,x}-y^{\mathrm{occ},-}_{r,x},
\qquad
\zeta_r=y^{\pi}_r\geq0,
\qquad
\omega^\star_{r,x}=y^w_{r,x,\gamma^\star_{r,x}}\geq0.
\]
For the fixed allocation orbit \(r\) and the selected grid points \(\gamma^\star_{r,x}\), dual feasibility gives the stationarity rows
\[
y^{\mathrm{norm},+}-y^{\mathrm{norm},-}
-\sum_x\alpha_{r,x}-\zeta_r=0
\]
and, for every compatible observation \(x\),
\[
\alpha_{r,x}-\omega^\star_{r,x}
+\sum_m\nu_{n,M,c}(m)P_m(x\mid r)
\{\gamma^\star_{r,x}-\tau_c(m)\}^2=0.
\]
Since the dual objective equals \(u=y^{\mathrm{norm},-}-y^{\mathrm{norm},+}\), these rows imply
\[
\begin{aligned}
u
&=-\sum_x\alpha_{r,x}-\zeta_r\\
&\leq \sum_x(-\alpha_{r,x})\\
&\leq
\sum_x\sum_m\nu_{n,M,c}(m)P_m(x\mid r)
\{\gamma^\star_{r,x}-\tau_c(m)\}^2\\
&=
\sum_m\nu_{n,M,c}(m)
\sum_x P_m(x\mid r)
\{\gamma^\star_{r,x}-\tau_c(m)\}^2.
\end{aligned}
\]
The posterior square decomposition gives
\[
\sum_m\nu_{n,M,c}(m)
\sum_x P_m(x\mid r)
\{\gamma^\star_{r,x}-\tau_c(m)\}^2
=
\operatorname{BayesRisk}_{n,c}(r,\nu_{n,M,c})
+
\sum_x D_{r,x,c}(\nu_{n,M,c})
\{\gamma^\star_{r,x}-\bar\tau_{r,x,c}(\nu_{n,M,c})\}^2.
\]
Indeed, for each \(x\) with positive predictive mass this is completion of the square around the posterior mean \(\bar\tau_{r,x,c}(\nu_{n,M,c})\); when the predictive mass is zero, all terms carrying \(P_m(x\mid r)\nu_{n,M,c}(m)\) vanish and the displayed convention contributes zero.  Moreover
\[
\sum_xD_{r,x,c}(\nu_{n,M,c})
=\sum_m\nu_{n,M,c}(m)\sum_xP_m(x\mid r)=1.
\]
Using the grid approximation bound for every \(x\),
\[
\sum_x D_{r,x,c}(\nu_{n,M,c})
\{\gamma^\star_{r,x}-\bar\tau_{r,x,c}(\nu_{n,M,c})\}^2
\leq
\frac{h_c^2}{4M^2}
=
\frac{C_0(c)}{4M^2}.
\]
Thus, for every allocation orbit \(r\),
\[
u
\leq
\operatorname{BayesRisk}_{n,c}(r,\nu_{n,M,c})
+
\frac{C_0(c)}{4M^2}.
\]
The allocation-orbit set is nonempty, so taking the infimum over \(r\) and using \(u=\Lambda_{n,M}(c)\) gives
\[
\Lambda_{n,M}(c)
\leq
B_{n,c}(\nu_{n,M,c})+\frac{C_0(c)}{4M^2}.
\]
Together with the preceding steps, this proves the full certificate sandwich. % lean: gridObjective_le_lowerCertificate_add_mesh

\item The orbit counts are stars-and-bars counts.  Response-count vectors are nonnegative counts over \(2^K\) binary response types with total \(n\), so
\[
|\mathcal M_{n,K}|=\binom{n+2^K-1}{2^K-1}.
\]
Allocation-count vectors are nonnegative counts over \(K\) arms with total \(n\), so
\[
|\mathcal R_{n,K}|=\binom{n+K-1}{K-1}.
\]
Finally, a pair \((r,x)\), with \(x\) compatible with \(r\), is equivalently a nonnegative count vector on arm-outcome cells \(\mathcal A_K\times\{0,1\}\) with total \(n\).  There are \(2K\) such cells, hence
\[
\sum_{r\in\mathcal R_{n,K}}|\{x:x\text{ is compatible with }r\}|
=
\binom{n+2K-1}{2K-1}.
\]
 % lean: allocation_observation_cardinality

\item Let \((a_k)\) be any positive sequence and let \(M_k\geq1\).  If
\[
\frac{a_k}{M_k^2}\longrightarrow0,
\]
then multiplication by the fixed constant \(C_0(c)/4\) gives
\[
a_k\,\frac{C_0(c)}{4M_k^2}
=
\frac{C_0(c)}{4}\,\frac{a_k}{M_k^2}
\longrightarrow0.
\]
For certificate sequences at mesh \(M_k\), the finite-sample sandwich gives, for all positive \(k\),
\[
B_k\leq\rho_k(c)\leq U_k,
\qquad
0\leq U_k-B_k\leq \frac{C_0(c)}{4M_k^2}.
\]
Finite initial values do not affect limits, and the preceding display therefore implies
\[
a_k(U_k-B_k)\longrightarrow0.
\]
Writing \(d_k(c)=C_0(c)/k-\rho_k(c)\), if \(a_kd_k(c)\to\ell\), then
\[
a_k\left\{\frac{C_0(c)}{k}-U_k\right\}
=
a_kd_k(c)-a_k\{U_k-\rho_k(c)\}
\longrightarrow\ell
\]
and
\[
a_k\left\{\frac{C_0(c)}{k}-B_k\right\}
=
a_kd_k(c)+a_k\{\rho_k(c)-B_k\}
\longrightarrow\ell,
\]
because both error terms are squeezed by \(a_k(U_k-B_k)\).  At the second-order scale \(a_k=k^{4/3}\) and with \(M_k=k\),
\[
\frac{k^{4/3}}{k^2}=k^{-2/3}\longrightarrow0,
\]
so
\[
k^{4/3}(U_k-B_k)\longrightarrow0.
\]
 % lean: rationalGridCertificateAsymptotics_proof
\end{enumerate}
\end{proof}

\begin{proof}[Proof of \cref{obj:thm:contrast-risk-continuity}]
\begin{enumerate}
\item First prove the one-sided estimate
\[
\sqrt{\rho_n(c)}\le \sqrt{\rho_n(c')}+\eta(c,c').
\]
Put
\[
\xi_a=c_a-c'_a,\qquad
\eta(c,c')=\frac12\sum_{a\in\mathcal A_K}|\xi_a|.
\]
For a binary response type \(t\), let
\[
\mathcal I(t)=\{a\in\mathcal A_K:t_a=1\}.
\]
Since both contrasts have zero sum, \(\sum_a \xi_a=0\), and therefore
\[
\sum_{a\in\mathcal I(t)}\xi_a
=
-\sum_{a\notin\mathcal I(t)}\xi_a .
\]
Thus
\[
2\left|\sum_{a\in\mathcal I(t)}\xi_a\right|
\le
\sum_{a\in\mathcal I(t)}|\xi_a|
+
\sum_{a\notin\mathcal I(t)}|\xi_a|
=
\sum_{a\in\mathcal A_K}|\xi_a|,
\]
so
\[
\left|\sum_{a\in\mathcal A_K}\xi_a t_a\right|
\le \eta(c,c').
\]
For any labeled schedule \(z\in\mathcal T_K^n\), using \(n\ge1\),
\[
\begin{aligned}
|\tau_c(z)-\tau_{c'}(z)|
&=
\left|\frac1n\sum_{i=1}^n\sum_{a\in\mathcal A_K}\xi_a z_{i,a}\right|  \\
&\le
\frac1n\sum_{i=1}^n
\left|\sum_{a\in\mathcal A_K}\xi_a z_{i,a}\right|
\le
\eta(c,c').
\end{aligned}
\]
Also set
\[
L_c=\sum_{a\in\mathcal A_K}|c_a|,
\qquad
I_c=\left[-\frac{L_c}{2},\frac{L_c}{2}\right].
\]
Applying the preceding target bound to the pair \(c,-c\) gives
\[
|2\tau_c(z)|=|\tau_c(z)-\tau_{-c}(z)|\le L_c,
\]
hence \(\tau_c(z)\in I_c\).

Fix any procedure \((\mathcal D,\widehat\tau')\) for contrast \(c'\). Define the transferred procedure for \(c\) by keeping the same assignment law and projecting the estimator output onto \(I_c\):
\[
\widehat\tau^{c}(A,Y)=
\operatorname{proj}_{I_c}\{\widehat\tau'(A,Y)\}.
\]
For \(y\in I_c\), projection onto the closed interval \(I_c\) satisfies
\[
\left|\operatorname{proj}_{I_c}(s)-y\right|^2\le |s-y|^2
\qquad(s\in\mathbb R),
\]
by the three cases \(s<-L_c/2\), \(s\in I_c\), and \(s>L_c/2\). For fixed \(z\) and assignment vector \(A\), define
\[
e_z(A)=\widehat\tau'(A,Y^{\mathrm{obs}}(z,A))-\tau_{c'}(z).
\]
Using \(\tau_c(z)\in I_c\) and the target bound above,
\[
\begin{aligned}
\left|\widehat\tau^{c}(A,Y^{\mathrm{obs}}(z,A))-\tau_c(z)\right|
&\le
\left|\widehat\tau'(A,Y^{\mathrm{obs}}(z,A))-\tau_c(z)\right|  \\
&\le
|e_z(A)|+|\tau_{c'}(z)-\tau_c(z)| \\
&\le
|e_z(A)|+\eta(c,c').
\end{aligned}
\]
For each target contrast \(\bar c\in\{c,c'\}\) and any estimator \(\widetilde\tau\), write its squared-error risk at the fixed schedule \(z\) as
\[
R_n^{\bar c}(\mathcal D,\widetilde\tau;z)
=
\mathbb E_{\mathcal D}\left[
\left\{
\widetilde\tau(A,Y^{\mathrm{obs}}(z,A))-\tau_{\bar c}(z)
\right\}^2
\right].
\]
Taking \(\mathcal D\)-expectations gives
\[
R_n^c(\mathcal D,\widehat\tau^c;z)
\le
\mathbb E_{\mathcal D}\left[(|e_z(A)|+\eta(c,c'))^2\right].
\]
Since \(\eta(c,c')\ge0\), Cauchy--Schwarz yields
\[
\mathbb E_{\mathcal D}|e_z(A)|
\le
\sqrt{\mathbb E_{\mathcal D}e_z(A)^2},
\]
and therefore
\[
\sqrt{\mathbb E_{\mathcal D}\left[(|e_z(A)|+\eta(c,c'))^2\right]}
\le
\sqrt{\mathbb E_{\mathcal D}e_z(A)^2}+\eta(c,c').
\]
Here \(\mathbb E_{\mathcal D}e_z(A)^2=R_n^{c'}(\mathcal D,\widehat\tau';z)\), so
\[
\sqrt{R_n^c(\mathcal D,\widehat\tau^c;z)}
\le
\sqrt{R_n^{c'}(\mathcal D,\widehat\tau';z)}+\eta(c,c').
\]
It remains to pass from this pointwise comparison to the two minimax values.  For each \(\bar c\in\{c,c'\}\), set
\[
\Psi_{\bar c}(\mathcal D,\widetilde\tau)
=
\max_{z\in\mathcal T_K^n} R_n^{\bar c}(\mathcal D,\widetilde\tau;z),
\]
so \cref{obj:def:labeled-schedule-game} says
\[
\rho_n(\bar c)=\inf_{\mathcal D,\widetilde\tau}\Psi_{\bar c}(\mathcal D,\widetilde\tau).
\]
The schedule set is finite and nonempty, and \(\sqrt{\cdot}\) is increasing on \([0,\infty)\), hence
\[
\max_{z\in\mathcal T_K^n}\sqrt{R_n^{\bar c}(\mathcal D,\widetilde\tau;z)}
=
\sqrt{\Psi_{\bar c}(\mathcal D,\widetilde\tau)}.
\]
We also have
\[
\sqrt{\rho_n(\bar c)}
=
\inf_{\mathcal D,\widetilde\tau}\sqrt{\Psi_{\bar c}(\mathcal D,\widetilde\tau)}.
\]
Indeed, every \(\Psi_{\bar c}(\mathcal D,\widetilde\tau)\) is at least \(\rho_n(\bar c)\), giving the \(\ge\) inequality. Conversely, the procedure class is nonempty, and for every \(\varepsilon>0\) there is a procedure with
\[
\Psi_{\bar c}(\mathcal D,\widetilde\tau)<\rho_n(\bar c)+\varepsilon;
\]
therefore the infimum of the square roots is at most \(\sqrt{\rho_n(\bar c)+\varepsilon}\), and continuity at \(\rho_n(\bar c)\) gives the reverse inequality as \(\varepsilon\downarrow0\).

Taking the maximum over \(z\) in the previous pointwise display gives
\[
\sqrt{\Psi_c(\mathcal D,\widehat\tau^c)}
\le
\sqrt{\Psi_{c'}(\mathcal D,\widehat\tau')}+\eta(c,c').
\]
Since \((\mathcal D,\widehat\tau^c)\) is one admissible procedure for contrast \(c\), taking the infimum over all procedures \((\mathcal D,\widehat\tau')\) for \(c'\) gives
\[
\sqrt{\rho_n(c)}\le \sqrt{\rho_n(c')}+\eta(c,c').
\]
% lean: root_minimax_one_sided_contrast_le

\item The contrast distance is symmetric:
\[
\eta(c',c)
=
\frac12\sum_{a\in\mathcal A_K}|c'_a-c_a|
=
\frac12\sum_{a\in\mathcal A_K}|c_a-c'_a|
=
\eta(c,c').
\]
Applying the one-sided estimate to the ordered pair \((c',c)\) gives
\[
\sqrt{\rho_n(c')}\le \sqrt{\rho_n(c)}+\eta(c',c)
=
\sqrt{\rho_n(c)}+\eta(c,c'),
\]
hence
\[
-\eta(c,c')\le \sqrt{\rho_n(c)}-\sqrt{\rho_n(c')}.
\]
% lean: contrast_risk_continuity

\item Applying the same one-sided estimate to the ordered pair \((c,c')\) gives
\[
\sqrt{\rho_n(c)}-\sqrt{\rho_n(c')}\le \eta(c,c').
\]
Together with the previous displayed lower bound, this is exactly
\[
\left|\sqrt{\rho_n(c)}-\sqrt{\rho_n(c')}\right|
\le
\eta(c,c').
\]
% lean: contrast_risk_continuity
\end{enumerate}
\end{proof}

\begin{proof}[Proof of \cref{obj:thm:real-contrast-grid-certificate-transfer}]
\begin{enumerate}
\item\label{it:real-contrast-finite-grid-transfer} Fix a rational zero-sum contrast \(q\in\mathbb Q^K\setminus\{0\}\) and integers \(n,M\ge 1\).  Apply \cref{obj:thm:rational-contrast-grid-certificate-sandwich} to \(q\).  It gives an invariant allocation certificate \(\pi\), a rational grid weight \(w\), a rational epigraph value \(u\), a rational response-count multiplier \(\nu\), and a barycenter rule \(\delta\), with
\[
B:=B_{n,q}(\nu),\qquad
U:=U_{n,M}(q;\pi,\delta),\qquad
u=\Lambda_{n,M}(q),
\]
and
\[
B\le \rho_n(q)\le U\le u
\le B+\frac{C_0(q)}{4M^2}.
\]
The multiplier \(\nu\) is a response-count probability prior.  Indeed, by \cref{obj:thm:rational-contrast-grid-certificate-sandwich}, \(\nu\) is the risk component of a dual-feasible multiplier for the grid-action program of \cref{obj:def:rational-contrast-grid-lp}.  In that program the response-count constraints are the inequality rows
\[
\sum_{r,x,g}P_m(x\mid r)\,w_{r,x,g}\{g-\tau_q(m)\}^2-u\le 0,
\qquad m\in\mathcal M_{n,K},
\]
so dual feasibility gives \(\nu(m)\ge0\) for every \(m\).  The epigraph variable \(u\) has objective coefficient \(1\), coefficient \(-1\) in each of those rows, and no other occurrence, so its dual stationarity equation is
\[
1-\sum_{m\in\mathcal M_{n,K}}\nu(m)=0.
\]
Thus
\[
\nu(m)\ge0\quad\text{for every }m\in\mathcal M_{n,K},
\qquad
\sum_{m\in\mathcal M_{n,K}}\nu(m)=1.
\]
With \(\nu\) now a probability prior, the lower certificate is an infimum of Bayes risks over the finite allocation-count set, hence \(B\ge0\); the displayed sandwich gives \(U\ge0\).  Put
\[
\eta_q:=\eta(c,q),\qquad
R^-:=\left\{\max(0,\sqrt B-\eta_q)\right\}^2,\qquad
R^+:=(\sqrt U+\eta_q)^2 .
\]

The pointwise transfer estimate used for both certificate directions is the following.  If a procedure for one contrast is clipped to the natural interval of another contrast, then for every schedule \(z\),
\[
\sqrt{R_n(\mathcal D,\widehat\tau^{\,c};z;c)}
\le
\sqrt{R_n(\mathcal D,\widehat\tau^{\,q};z;q)}+\eta_q .
\]
Indeed, for each assignment \(A\),
\[
|\tau_c(z)-\tau_q(z)|\le \eta_q
\]
because the target difference is an average of binary response-type scores and
\[
\eta_q=\frac12\sum_{a\in\mathcal A_K}|c(a)-q(a)|.
\]
Projection onto \([-L_c/2,L_c/2]\) can only reduce squared distance to \(\tau_c(z)\), and Cauchy--Schwarz gives
\[
\mathbb E_{\mathcal D}\!\left(|E_A|+\eta_q\right)^2
\le
\left\{\sqrt{\mathbb E_{\mathcal D}E_A^2}+\eta_q\right\}^2
\]
for the original residual \(E_A=\widehat\tau^{\,q}(A,Y^{\mathrm{obs}})-\tau_q(z)\).

The upper procedure keeps the invariant allocation law induced by \(\pi\) and uses
\[
\widehat\tau^{\,c}(A,Y^{\mathrm{obs}})
=
\operatorname{proj}_{[-L_c/2,L_c/2]}
\{\delta(r(A),x(A,Y^{\mathrm{obs}}))\}.
\]
The rational certificate bounds the \(q\)-risk of the barycenter rule by \(U\), so the transfer estimate yields worst-case \(c\)-risk at most \(R^+\), and therefore
\[
\rho_n(c)\le R^+ .
\]

For the lower endpoint, spread this response-count prior uniformly over each labeled schedule orbit, and write the resulting expectation as
\[
\mathbb E^{\mathrm{sch}}_{\nu}[\varphi]
:=
\sum_{z\in\mathcal T_K^n}
\frac{\nu(m(z))}{|\{z':m(z')=m(z)\}|}\,\varphi(z).
\]
The denominator is positive because each response-count vector is represented by at least one labeled schedule, and the displayed nonnegativity and mass identity make \(\mathbb E^{\mathrm{sch}}_{\nu}\) a probability expectation.  We first record the all-procedure lower bound supplied by this lifted prior.  Let \(\mathfrak p=(\mathcal D_{\mathfrak p},\widehat\tau_{\mathfrak p})\) be any labeled procedure for contrast \(q\), and average \(\mathfrak p\) over the unit-permutation group \(\mathfrak S_n\).  By the lossless symmetrization and exact invariant correspondence in \cref{obj:thm:exact-response-type-game}, the averaged procedure is invariant and corresponds to an orbit procedure \(\mathfrak o\).  If \(\alpha_{\mathfrak o}\) is its allocation-count law, define its orbit risk at response-count vector \(m\) by
\[
\operatorname{Risk}^{\mathrm{orb}}_{n,q}(\mathfrak o;m)
:=
\sum_{r\in\mathcal R_{n,K}}\alpha_{\mathfrak o}(r)
\sum_x P_m(x\mid r)
\{\widehat\tau_{\mathfrak o}(r,x)-\tau_q(m)\}^2 .
\]
For a fixed allocation count \(r\), the conditional Bayes risk under \(\nu\) is bounded below by \(B=B_{n,q}(\nu)\); on zero prior-predictive fibers the corresponding summand is zero, and on positive prior-predictive fibers the posterior-mean residual is the conditional minimum.  Hence, for every orbit procedure,
\[
B
\le
\sum_{m}\nu(m)\operatorname{Risk}^{\mathrm{orb}}_{n,q}(\mathfrak o;m).
\]
Because the lifted prior is uniform within each response-count orbit,
\[
\sum_m\nu(m)\operatorname{Risk}^{\mathrm{orb}}_{n,q}(\mathfrak o;m)
=
\mathbb E^{\mathrm{sch}}_{\nu}
\left[\operatorname{Risk}^{\mathrm{orb}}_{n,q}(\mathfrak o;m(z))\right].
\]
Lossless symmetrization gives, for every labeled schedule \(z\),
\[
\operatorname{Risk}^{\mathrm{orb}}_{n,q}(\mathfrak o;m(z))
\le
\frac1{|\mathfrak S_n|}
\sum_{\sigma\in\mathfrak S_n}
R_n(\mathcal D_{\mathfrak p},\widehat\tau_{\mathfrak p};\sigma z;q).
\]
The lifted prior is invariant under unit permutations, so averaging the last display over \(z\) yields
\[
B
\le
\mathbb E^{\mathrm{sch}}_{\nu}
\left[R_n(\mathcal D_{\mathfrak p},\widehat\tau_{\mathfrak p};z;q)\right]
\]
for every labeled \(q\)-procedure \(\mathfrak p\).

Now let \(\mathfrak r=(\mathcal D_{\mathfrak r},\widehat\tau_{\mathfrak r})\) be any procedure for the real contrast \(c\), and let \(\mathfrak r^{q}\) be the procedure obtained by clipping its action to \([-L_q/2,L_q/2]\).  Applying the pointwise transfer estimate in the direction from \(c\) to \(q\),
\[
\sqrt{R_n(\mathfrak r^{q};z;q)}
\le
\sqrt{R_n(\mathcal D_{\mathfrak r},\widehat\tau_{\mathfrak r};z;c)}+\eta_q .
\]
Set
\[
\mathscr V_c(\mathfrak r)
:=
\mathbb E^{\mathrm{sch}}_{\nu}
\left[R_n(\mathcal D_{\mathfrak r},\widehat\tau_{\mathfrak r};z;c)\right].
\]
Squaring the preceding inequality, integrating with respect to \(\mathbb E^{\mathrm{sch}}_{\nu}\), and using Cauchy--Schwarz in the finite schedule space gives
\[
\mathbb E^{\mathrm{sch}}_{\nu}
\left[R_n(\mathfrak r^{q};z;q)\right]
\le
\mathscr V_c(\mathfrak r)
+2\eta_q\sqrt{\mathscr V_c(\mathfrak r)}
+\eta_q^2
=
\{\sqrt{\mathscr V_c(\mathfrak r)}+\eta_q\}^2.
\]
Combining this with the all-procedure \(q\)-lower bound applied to \(\mathfrak r^q\) gives
\[
B\le\{\sqrt{\mathscr V_c(\mathfrak r)}+\eta_q\}^2.
\]
If \(B\le0\), then \(\sqrt B=0\) and \(R^-=0\le \mathscr V_c(\mathfrak r)\).  If \(B>0\), the displayed inequality gives \(\sqrt B\le\sqrt{\mathscr V_c(\mathfrak r)}+\eta_q\).  When \(\sqrt B\le\eta_q\) this again gives \(R^-=0\le\mathscr V_c(\mathfrak r)\), and otherwise
\[
R^-=(\sqrt B-\eta_q)^2\le \mathscr V_c(\mathfrak r).
\]
Thus the lifted schedule prior has Bayes risk at least \(R^-\) against every real-contrast procedure.  Since prior risk is bounded above by worst-case risk for each procedure, taking the minimax value gives
\[
R^-\le \rho_n(c).
\]

It remains to verify the stated width.  From the rational sandwich,
\[
U-B\le \frac{C_0(q)}{4M^2}.
\]
Also \cref{obj:thm:embedded-two-arm-converse} gives
\[
\rho_n(q)\le \frac{C_0(q)}{n},
\]
so
\[
U\le \frac{C_0(q)}{n}+\frac{C_0(q)}{4M^2}.
\]
Consequently,
\[
\sqrt U
\le
\sqrt{C_0(q)}\,n^{-1/2}
+
\frac{\sqrt{C_0(q)}}{2M},
\]
because the square of the right-hand side contains the displayed upper bound plus a nonnegative cross term.  Finally,
\[
B-\{\max(0,\sqrt B-\eta_q)\}^2
\le
2\eta_q\sqrt B+\eta_q^2,
\]
and \(\sqrt B\le\sqrt U\), whence
\[
R^+-R^-
\le
\frac{C_0(q)}{4M^2}+4\eta_q\sqrt U+2\eta_q^2.
\]
Substituting the bound on \(\sqrt U\) and using
\[
\frac{2\sqrt{C_0(q)}\,\eta_q}{M}
\le
\frac{C_0(q)}{M^2}+\eta_q^2
\]
gives
\[
R^+-R^-\le
\frac{5C_0(q)}{4M^2}
+4\sqrt{C_0(q)}\,\eta_q\,n^{-1/2}
+3\eta_q^2.
\]
Thus \(B,U,R^-,R^+\) form a transferred certificate for \((n,M,c,q)\), proving the finite-grid transfer assertion. % lean: fixedRealContrastTransferCertificate

\item We next construct the rational approximating sequence.  Since \(c\) is nonzero, \(L_c>0\).  For \(s\in\mathbb N\), define the approximation radius
\[
\operatorname{rad}_s=
\begin{cases}
L_c/4,& s=0,\\
\min\{L_c/4,s^{-3}\},& s\ge1 .
\end{cases}
\]
Then \(0<\operatorname{rad}_s<L_c/2\).  Choose a fixed arm \(a_\circ\).  By density of \(\mathbb Q\) in \(\mathbb R\), choose rational numbers \(\widetilde c_s(a)\) for \(a\ne a_\circ\) with
\[
|c(a)-\widetilde c_s(a)|<\frac{\operatorname{rad}_s}{2K},
\]
and set
\[
\widetilde c_s(a_\circ)
=
-\sum_{a\ne a_\circ}\widetilde c_s(a).
\]
Because \(c\) is zero-sum,
\[
c(a_\circ)=-\sum_{a\ne a_\circ}c(a),
\]
and hence
\[
|c(a_\circ)-\widetilde c_s(a_\circ)|
\le
\sum_{a\ne a_\circ}|c(a)-\widetilde c_s(a)|.
\]
Therefore
\[
\sum_{a\in\mathcal A_K}|c(a)-\widetilde c_s(a)|
\le \operatorname{rad}_s .
\]
Since \(\operatorname{rad}_s<L_c/2\), this rational zero-sum vector is nonzero; otherwise \(L_c\le\operatorname{rad}_s\).  Denote the resulting rational contrast by \(c^{(s)}\).  Then
\[
\eta(c,c^{(s)})\le \operatorname{rad}_s,
\]
and for \(s\ge1\),
\[
\eta(c,c^{(s)})\le s^{-3}.
\]
It follows that
\[
\eta(c,c^{(s)})\to0,
\qquad
s^{5/6}\eta(c,c^{(s)})
\le s^{5/6}s^{-3}=s^{-13/6}\to0 .
\]
For each arm \(a\),
\[
|c^{(s)}(a)-c(a)|
\le
\sum_{b\in\mathcal A_K}|c^{(s)}(b)-c(b)|
=
2\eta(c,c^{(s)}),
\]
so \(c^{(s)}(a)\to c(a)\). % lean: real_contrast_grid_certificate_transfer

\item For each \(s\ge1\), apply the fixed construction from \cref{it:real-contrast-finite-grid-transfer} with \(q=c^{(s)}\), \(n=s\), and \(M=s\), obtaining \(B_s,U_s,R_s^-,R_s^+\).  The construction gives a transferred certificate for every \(s\ge1\).  Write
\[
\eta_s:=\eta(c,c^{(s)}),
\qquad
C_s:=C_0(c^{(s)}).
\]
Because \(c^{(s)}(a)\to c(a)\) for every arm and the arm set is finite,
\[
L_{c^{(s)}}=\sum_{a\in\mathcal A_K}|c^{(s)}(a)|
\longrightarrow
\sum_{a\in\mathcal A_K}|c(a)|=L_c,
\]
and hence
\[
C_s=\frac{L_{c^{(s)}}^2}{4}\longrightarrow C_0(c).
\]
The certificate width bound gives, for \(s\ge1\),
\[
0\le s^{4/3}(R_s^+-R_s^-)
\le
s^{4/3}\left\{
\frac{5C_s}{4s^2}
+4\sqrt{C_s}\,\eta_s\,s^{-1/2}
+3\eta_s^2
\right\}.
\]
The three terms on the right converge to zero:
\[
s^{4/3}\frac{5C_s}{4s^2}
=
\frac{5C_s}{4}s^{-2/3}\to0,
\]
\[
s^{4/3}\,4\sqrt{C_s}\,\eta_s\,s^{-1/2}
=
4\sqrt{C_s}\,\{s^{5/6}\eta_s\}\to0,
\]
and, using \(\eta_s\le s^{-3}\),
\[
s^{4/3}\,3\eta_s^2
\le
3s^{4/3}s^{-6}
=
3s^{-14/3}\to0.
\]
By squeezing,
\[
s^{4/3}(R_s^+-R_s^-)\to0.
\]
This proves the real-contrast approximation assertion. % lean: real_contrast_grid_certificate_transfer

\item Finally suppose a rational zero-sum contrast \(q\) embeds coefficientwise as \(c\).  For any \(n\ge1\), apply the fixed construction with \(M=n\).  The rational sandwich gives
\[
U-B\le \frac{C_0(q)}{4n^2}.
\]
Since the real embedding of \(q\) is \(c\), \(C_0(q)=C_0(c)\), and hence
\[
U-B\le \frac{C_0(c)}{4n^2}.
\]
Together with the transferred certificate constructed in \cref{it:real-contrast-finite-grid-transfer}, this proves the exact rational case. % lean: real_contrast_grid_certificate_transfer
\end{enumerate}
\end{proof}

\begin{proof}[Proof of \cref{obj:thm:k3-grid-certificate-sandwich}]
Fix \(n\geq 1\) and \(M\geq 1\).
\begin{enumerate}
\item Let
\[
c_{\mathbb Q}^\dagger=(1,-1/2,-1/2)\in\mathbb Q^3,
\qquad
c^\dagger=(1,-1/2,-1/2)\in\mathbb R^3 .
\]
This contrast is nonzero and zero-sum, and
\[
L_{c^\dagger}=|1|+|-1/2|+|-1/2|=2,
\qquad
C_0(c^\dagger)=\left(\frac{L_{c^\dagger}}2\right)^2=1 .
\]
Hence the real embedding of \(c_{\mathbb Q}^\dagger\) has \(C_0=1\), and the grid in \cref{obj:def:rational-contrast-grid-lp} becomes
\[
\Gamma_{M,c^\dagger}
=
\{-1+j/M:0\leq j\leq 2M\}
=
\Gamma_M .
\]
Therefore, after specializing \cref{obj:thm:rational-contrast-grid-certificate-sandwich} to \(K=3\) and \(c=c_{\mathbb Q}^\dagger\), the rational contrast grid program of \cref{obj:def:rational-contrast-grid-lp} is the three-arm rational grid program of \cref{obj:def:k3-rational-lp}, with
\[
\Lambda_{n,M}(c^\dagger)=\Lambda_{n,M}^\dagger,
\qquad
\rho_n(c^\dagger)=\rho_n^\dagger .
\]
% lean: k3_grid_certificate_sandwich

\item The specialized certificate theorem supplies a rational allocation certificate \(\pi^{n,M}\), a joint design--action weight \(w^{n,M}\), a rational epigraph value \(u\), a rational response-count multiplier \(\nu_{n,M}\), and an estimator \(\delta_{n,M}\). Its exact-certificate conclusion gives the rational grid weight whose real embedding is \(w^{n,M}\), primal feasibility, dual feasibility, the identification of the risk component with \(\nu_{n,M}\), and dual objective value \(u\). Its barycenter conclusion gives, for every \(r\in\mathcal R_{n,3}\) and compatible \(x\),
\[
\delta_{n,M}(r,x)
=
\begin{cases}
\displaystyle
\sum_{g\in\Gamma_M}g\,\frac{w_{r,x,g}^{n,M}}{\pi_r^{n,M}},
& \pi_r^{n,M}>0,\\[1ex]
0,& \pi_r^{n,M}=0.
\end{cases}
\]
The second branch fixes the estimator at allocation counts that carry no certificate mass; those counts are never charged by the certificate, so the choice is immaterial.

The same specialization gives
\[
u=\Lambda_{n,M}^\dagger .
\]
For a response-count multiplier \(\nu\), the specialized lower endpoint is
\[
B_{n,c^\dagger}(\nu)
=
\inf_{r\in\mathcal R_{n,3}}
\operatorname{BayesRisk}_{n,c^\dagger}(r,\nu)
=
B_n(\nu),
\]
and the specialized upper endpoint is precisely
\[
U_{n,M}
=
\sup_{m\in\mathcal M_{n,3}}
\sum_{r\in\mathcal R_{n,3}}\pi_r^{n,M}
\sum_x
P_m(x\mid r)
\{\delta_{n,M}(r,x)-\tau_{c^\dagger}(m)\}^2 .
\]
Thus the certificate sandwich in \cref{obj:thm:rational-contrast-grid-certificate-sandwich}, together with \(C_0(c^\dagger)=1\), yields
\[
B_n(\nu_{n,M})
\le
\rho_n^\dagger
\le
U_{n,M}
\le
\Lambda_{n,M}^\dagger
\le
B_n(\nu_{n,M})+\frac{1}{4M^2}.
\]
% lean: rational_contrast_grid_certificate_sandwich

\item Let \(a_k>0\) be any positive normalizing sequence and let \(M_k\geq1\) for all \(k\). If
\[
\frac{a_k}{M_k^2}\longrightarrow 0,
\]
then
\[
a_k\,\frac{1}{4M_k^2}
=
\frac14\,\frac{a_k}{M_k^2}
\longrightarrow 0 .
\]
This is the mesh-rate conclusion at \(C_0(c^\dagger)=1\). % lean: k3_grid_certificate_sandwich

\item Now let \(a_k>0\), \(M_k\), \(L_k\), \(U_k\), and \(C\) satisfy the hypotheses of the final assertion. All inequalities below are for \(k\geq1\), which is an eventual range for limits along \(k\to\infty\). From
\[
L_k\leq \rho_k^\dagger\leq U_k,
\qquad
U_k-L_k\leq \frac{1}{4M_k^2},
\]
we have
\[
0\leq a_k(\rho_k^\dagger-L_k)
\leq
a_k(U_k-L_k)
\leq
a_k\,\frac{1}{4M_k^2},
\]
and also
\[
0\leq a_k(U_k-\rho_k^\dagger)
\leq
a_k(U_k-L_k)
\leq
a_k\,\frac{1}{4M_k^2}.
\]
The mesh assumption and the preceding step give
\[
a_k(\rho_k^\dagger-L_k)\longrightarrow0,
\qquad
a_k(U_k-\rho_k^\dagger)\longrightarrow0 .
\]
Finally,
\[
d_k(c^\dagger)
=
\frac{C_0(c^\dagger)}{k}-\rho_k^\dagger .
\]
Therefore
\[
a_k\left\{\frac{C_0(c^\dagger)}{k}-L_k\right\}
=
a_k d_k(c^\dagger)
+
a_k(\rho_k^\dagger-L_k),
\]
and
\[
a_k\left\{\frac{C_0(c^\dagger)}{k}-U_k\right\}
=
a_k d_k(c^\dagger)
-
a_k(U_k-\rho_k^\dagger).
\]
Since \(a_kd_k(c^\dagger)\to C\) and both residual terms converge to \(0\), the two displayed sequences converge to \(C\). % lean: k3_grid_certificate_sandwich
\end{enumerate}
\end{proof}

\begin{proof}[Proof of \cref{obj:prop:k3-scalar-score-not-minimax-preserving}]
\begin{enumerate}
\item
For \(c^\dagger=(1,-1/2,-1/2)\), the contrast-weighted one-unit law is
\[
\Pr(A_i=1)=\frac{|1|}{|1|+|-1/2|+|-1/2|}=\frac12,
\qquad
\Pr(A_i=2)=\Pr(A_i=3)=\frac14 .
\]
Since \(c^\dagger\) is positive on arm \(1\) and negative on arms \(2,3\),
\[
Z_i=
\begin{cases}
 2t_{i,1}-1, & A_i=1,\\
 1-2t_{i,2}, & A_i=2,\\
 1-2t_{i,3}, & A_i=3.
\end{cases}
\]
Thus, conditionally on \(z\),
\[
\begin{aligned}
\mathbb E_{\mathcal D^\star}(Z_i\mid z)
&=\frac12(2t_{i,1}-1)+\frac14(1-2t_{i,2})
  +\frac14(1-2t_{i,3})  \\
&=t_{i,1}-\frac{t_{i,2}+t_{i,3}}2
=\mu_i .
\end{aligned}
\]
This proves the asserted conditional mean for every \(n\ge1\), schedule \(z\), and unit \(i\). % lean: k3SignedScore_mean

\item
The score is always sign-valued, so \(Z_i^2=1\). Combining this with the preceding mean identity gives
\[
\operatorname{Var}_{\mathcal D^\star}(Z_i\mid z)
=\mathbb E_{\mathcal D^\star}(Z_i^2\mid z)
-\{\mathbb E_{\mathcal D^\star}(Z_i\mid z)\}^2
=1-\mu_i^2 .
\]
% lean: k3SignedScore_var

\item
Write
\[
\overline Z=\frac1n\sum_i Z_i .
\]
The assignment law is a product law, so the \(Z_i\)'s are conditionally independent given \(z\). Moreover
\[
\tau_{c^\dagger}(z)
=\frac1n\sum_i\left(t_{i,1}-\frac{t_{i,2}+t_{i,3}}2\right)
=\frac1n\sum_i\mu_i ,
\]
and therefore \(\overline Z\) is conditionally unbiased for \(\tau_{c^\dagger}(z)\). Hence, for \(n\ge1\),
\[
\begin{aligned}
\mathbb E_{\mathcal D^\star}\!\left[
\{\overline Z-\tau_{c^\dagger}(z)\}^2\,\middle|\,z\right]
&=\operatorname{Var}_{\mathcal D^\star}(\overline Z\mid z)  \\
&=\frac1{n^2}\sum_i \operatorname{Var}_{\mathcal D^\star}(Z_i\mid z)\\
&=\frac1{n^2}\sum_i(1-\mu_i^2)
=\frac1n-\frac1{n^2}\sum_i\mu_i^2 .
\end{aligned}
\]
% lean: k3SignedScore_average_mse

\item
For the scalar experiment at \(n=3\), let
\[
\mathfrak L_Z=\{-1,-1/2,0,1/2,1\}
\]
be the feasible set of scalar means. Given \(\boldsymbol\mu=(\mu_1,\mu_2,\mu_3)\in\mathfrak L_Z^3\), the retained observation has independent coordinates with
\[
\Pr_\mu(Z_i=1)=\frac{1+\mu_i}{2},
\qquad
\Pr_\mu(Z_i=-1)=\frac{1-\mu_i}{2},
\]
and the scalar target is \(M_\mu/3\), where
\[
M_\mu=\sum_{i=1}^3\mu_i,
\qquad
Q_\mu=\sum_{i=1}^3\mu_i^2 .
\]
Set
\[
\alpha=\frac{3-\sqrt3}{6},
\qquad
\widehat\theta_Z(Z_1,Z_2,Z_3)=\alpha\sum_{i=1}^3 Z_i .
\]
Then
\[
\begin{aligned}
\mathbb E_\mu\!\left[
\left\{\widehat\theta_Z-\frac{M_\mu}{3}\right\}^2\right]
&=\alpha^2(3-Q_\mu)+\left(\alpha-\frac13\right)^2M_\mu^2  \\
&=\frac{2-\sqrt3}{18}\{9-3Q_\mu+M_\mu^2\}.
\end{aligned}
\]
By Cauchy's inequality, \(M_\mu^2\le3Q_\mu\), and hence
\[
\mathbb E_\mu\!\left[
\left\{\widehat\theta_Z-\frac{M_\mu}{3}\right\}^2\right]
\le \frac{2-\sqrt3}{2}
=1-\frac{\sqrt3}{2}.
\]
Taking the supremum over \(\boldsymbol\mu\in\mathfrak L_Z^3\) and then the infimum over scalar rules gives
\[
V_3^Z\le 1-\frac{\sqrt3}{2}.
\]
% lean: k3ScalarMinimaxValue_le_exact

\item
For the reverse inequality, put the following weights on the five homogeneous states
\[
\lambda_j=-1+\frac j2,\qquad j=0,1,2,3,4:
\]
\[
\varpi_0=\varpi_4=\frac{-11+7\sqrt3}{12},\qquad
\varpi_1=\varpi_3=\frac{8-4\sqrt3}{3},\qquad
\varpi_2=\frac{-5+3\sqrt3}{2}.
\]
They are nonnegative and satisfy \(\sum_{j=0}^4\varpi_j=1\). For any scalar rule \(\widehat\theta\), its Bayes risk under this prior is
\[
\mathcal B(\widehat\theta)
=\sum_{j=0}^4\varpi_j\,
\mathbb E_{\lambda_j}\!\left[(\widehat\theta(Z)-\lambda_j)^2\right],
\]
where \(\mathbb E_{\lambda_j}\) denotes the product law with all three scalar means equal to \(\lambda_j\). If \(N_+=\#\{i:Z_i=1\}\), then for \(b=0,1,2,3\) the prior predictive mass of the event \(\{N_+=b\}\) and its target numerator are
\[
\mathfrak P_b
=\binom{3}{b}\sum_{j=0}^4\varpi_j
\left(\frac{1+\lambda_j}{2}\right)^b
\left(\frac{1-\lambda_j}{2}\right)^{3-b},
\qquad
\mathfrak H_b
=\binom{3}{b}\sum_{j=0}^4\varpi_j\lambda_j
\left(\frac{1+\lambda_j}{2}\right)^b
\left(\frac{1-\lambda_j}{2}\right)^{3-b}.
\]
A direct expansion gives
\[
\begin{array}{c|c|c|c}
b&\mathfrak P_b&\mathfrak H_b&\mathfrak H_b/\mathfrak P_b\\ \hline
0&(-1+3\sqrt3)/16&(6-5\sqrt3)/16&-3\alpha\\
1&(9-3\sqrt3)/16&(-6+3\sqrt3)/16&-\alpha\\
2&(9-3\sqrt3)/16&(6-3\sqrt3)/16&\alpha\\
3&(-1+3\sqrt3)/16&(-6+5\sqrt3)/16&3\alpha .
\end{array}
\]
Thus the posterior mean is exactly \(\alpha(2N_+-3)=\widehat\theta_Z(Z)\), and completing the square yields
\[
\mathcal B(\widehat\theta)
=
\mathcal B(\widehat\theta_Z)
+\sum_{Z\in\{-1,1\}^3}
\Pr_{\mathrm{prior}}(Z)
\{\widehat\theta(Z)-\widehat\theta_Z(Z)\}^2
\ge \mathcal B(\widehat\theta_Z).
\]
For a homogeneous level \(\lambda_j\),
\[
\mathbb E_{\lambda_j}\!\left[(\widehat\theta_Z(Z)-\lambda_j)^2\right]
=
\alpha^2\,3(1-\lambda_j^2)+(3\alpha\lambda_j-\lambda_j)^2,
\]
and summing this identity with the weights \(\varpi_j\) gives
\[
\mathcal B(\widehat\theta_Z)=1-\frac{\sqrt3}{2}.
\]
Every supremum risk dominates this Bayes risk, so
\[
1-\frac{\sqrt3}{2}\le V_3^Z .
\]
Together with the preceding item, this proves \(V_3^Z=1-\sqrt3/2\). % lean: k3ScalarMinimaxValue_ge_exact

\item
The full-data rule is defined from the observed arm labels and outcomes by
\[
r_1=\sum_i\mathbf 1\{A_i=1\},\qquad
s_1=\sum_{i:A_i=1}(2Y_i^{\mathrm{obs}}-1),\qquad
s_-=\sum_{i:A_i\in\{2,3\}}(1-2Y_i^{\mathrm{obs}}).
\]
The rule is given by a lookup table.  Let \(\mathrm T_F(r_1,s_1,s_-)\) be the integer in the table
\[
\begin{array}{c|r@{\qquad}c|r}
(0,0,-3)&-623&(0,0,-1)&-183\\
(0,0,1)&183&(0,0,3)&623\\
(1,-1,-2)&-638&(1,-1,0)&-199\\
(1,-1,2)&148&(1,1,-2)&-148\\
(1,1,0)&199&(1,1,2)&638\\
(2,-2,-1)&-648&(2,-2,1)&-210\\
(2,0,-1)&-167&(2,0,1)&167\\
(2,2,-1)&210&(2,2,1)&648\\
(3,-3,0)&-660&(3,-1,0)&-167\\
(3,1,0)&167&(3,3,0)&660 .
\end{array}
\]
The displayed triples are exactly the triples \((r_1,s_1,s_-)\) that assignments and binary outcomes can generate, so the rule is defined wherever it is used; set \(\mathrm T_F=0\) elsewhere. Define
\[
\widehat\tau_F(A,Y^{\mathrm{obs}})
=
\operatorname{proj}_{[-1,1]}\!\left(\frac{\mathrm T_F(r_1,s_1,s_-)}{1000}\right).
\]
All displayed values lie in \([-1000,1000]\), so the projection leaves the table value divided by \(1000\) unchanged. For any schedule \(z\), its risk under \(\mathcal D^\star\) is the finite rational sum
\[
R_F(z)
=
\sum_{a\in\mathcal A_3^3}
\prod_{i=1}^3 q^\star_{a_i}
\left\{
\frac{\mathrm T_F(r_1(a),s_1(a,z),s_-(a,z))}{1000}
-\tau_{c^\dagger}(z)
\right\}^2 .
\]
Identifying an assignment \(a\in\mathcal A_3^3\) with its ordered triple gives a \(27\)-term rational sum. To certify the uniform bound, write each schedule as
\[
z=(\mathbf t^{(1)},\mathbf t^{(2)},\mathbf t^{(3)}),
\qquad
\mathbf t^{(j)}\in\{0,1\}^3,
\]
and order the eight possible response types as
\[
\mathfrak O_3=(000,001,010,011,100,101,110,111).
\]
For \(\mathbf t^{(1)},\mathbf t^{(2)},\mathbf t^{(3)}\in\mathfrak O_3\), define the cleared slack
\[
\Xi(\mathbf t^{(1)},\mathbf t^{(2)},\mathbf t^{(3)})
=
144000000\left\{\frac{511653}{4000000}
-R_F(\mathbf t^{(1)},\mathbf t^{(2)},\mathbf t^{(3)})\right\}.
\]
Substituting the \(27\)-term expression for \(R_F\) and clearing denominators makes \(\Xi\) an integer-valued function of the triple, and the check is a finite one: there are \(8^3=512\) triples, and \(\Xi\ge 0\) at every one of them.  The evaluation is an exact rational computation: after substituting the displayed formula for \(R_F\), the \(512\) integer values of \(\Xi\) are all nonnegative.  The bound is attained: for instance
\[
\Xi(000,111,111)=0,
\]
so no smaller constant than \(511653/4000000\) works for this estimator.
Since \(\Xi\ge0\) on all of \((\{0,1\}^3)^3\),
\[
R_F(z)\le \frac{511653}{4000000}
\qquad\text{for every }z\in\mathcal T_3^3 .
\]
For the schedule with response types \((0,0,0),(1,1,1),(1,1,1)\), the target is \(0\), and the same finite sum reduces to
\[
\begin{aligned}
R_F(z)
&=\frac18\left(\frac{183}{1000}\right)^2
+\frac18\left(\frac{638}{1000}\right)^2
+\frac14\left(\frac{199}{1000}\right)^2
+\frac14\left(\frac{167}{1000}\right)^2  \\
&\qquad
+\frac18\left(\frac{648}{1000}\right)^2
+\frac18\left(\frac{167}{1000}\right)^2
=\frac{511653}{4000000}.
\end{aligned}
\]
Therefore the worst-case risk of this full-data rule is exactly \(511653/4000000\). % lean: k3FullDataRule_worstCaseRisk

\item
The rational comparison in the full-data certificate is
\[
\frac{511653}{4000000}<\frac{18213}{136000}.
\]
% lean: k3_rational_separation
Also,
\[
\sqrt3<\frac{117787}{68000},
\]
because squaring both positive sides gives \(3<(117787/68000)^2\). Hence
\[
1-\frac{\sqrt3}{2}
>
1-\frac1{2}\frac{117787}{68000}
=
\frac{18213}{136000}.
\]
Thus
\[
\frac{511653}{4000000}
<
\frac{18213}{136000}
<
1-\frac{\sqrt3}{2}.
\]
% lean: k3_scalar_score_not_minimax_preserving

\item
It remains to compute the scalar moment envelope. Let
\[
\mathfrak L_Z=\{-1,-1/2,0,1/2,1\},
\qquad
M=\left|\sum_i\mu_i\right|,
\qquad
Q_\mu=\sum_i\mu_i^2,
\]
with each \(\mu_i\in\mathfrak L_Z\). Pointwise on \(\mathfrak L_Z\),
\[
\frac{|\mu_i|}{2}\le \mu_i^2,
\qquad
\frac{3|\mu_i|-1}{2}\le \mu_i^2 .
\]
Since \(\sum_i|\mu_i|\ge M\), summing these two inequalities gives
\[
Q_\mu\ge \frac M2,
\qquad
Q_\mu\ge \frac{3M-n}{2}.
\]
Consequently
\[
Q_\mu\ge
\begin{cases}
M/2, & M\le n/2,\\
(3M-n)/2, & M>n/2.
\end{cases}
\]
For attainment, feasibility implies \(M=k/2\) for some integer \(0\le k\le2n\). If \(k\le n\), take \(k\) coordinates equal to \(1/2\) and the remaining coordinates equal to \(0\). Then
\[
\left|\sum_i\mu_i\right|=\frac k2=M,
\qquad
\sum_i\mu_i^2=\frac k4=\frac M2 .
\]
If \(k>n\), write \(\ell=2n-k\), so \(0\le\ell<n\), take \(\ell\) coordinates equal to \(1/2\), and take the remaining \(n-\ell\) coordinates equal to \(1\). Then
\[
\left|\sum_i\mu_i\right|
=n-\frac{\ell}{2}
=\frac k2=M,
\]
and
\[
\sum_i\mu_i^2
=n-\frac{3\ell}{4}
=\frac{3M-n}{2}.
\]
For \(n=0\), feasibility forces \(M=0\), and the empty vector gives value \(0\), matching the first branch. Taking the infimum over feasible scalar mean vectors yields
\[
\min\Bigl\{Q_\mu:\ \mu\in\mathfrak L_Z^{\,n},\ \Bigl|\sum_i\mu_i\Bigr|=M\Bigr\}
=
\begin{cases}
M/2, & M\le n/2,\\
(3M-n)/2, & M>n/2.
\end{cases}
\]
Equivalently, in the notation of the statement, this is the asserted formula for the minimum attainable \(\sum_i\mu_i^2\). % lean: scalarMomentMinimum_wedge

\item
Combining the score identities, the average-score risk formula, the exact scalar minimax value, the full-data witness with its rational separation, and the scalar moment envelope proves all asserted clauses. % lean: k3_scalar_score_not_minimax_preserving
\end{enumerate}
\end{proof}

\begin{proof}[Proof of \cref{obj:thm:second-order-rate-and-certificate-frontier}]
Let
\[
s_n=n^{4/3},
\qquad
x_n=x_n(c)=s_n d_n(c),
\qquad
d_n(c)=\frac{C_0(c)}{n}-\rho_n(c),
\]
Since \(C_0(c)=L_c^2/4\) and the contrast is nonzero, \(C_0(c)>0\).

\begin{enumerate}
\item \textbf{Eventual two-sided bounds.}
\Cref{obj:thm:universal-second-order-rate} supplies \(\kappa_c>0\), the liminf--limsup chain
\[
C_0(c)\kappa_c
\leq
\liminf_{n\to\infty}x_n
\leq
\limsup_{n\to\infty}x_n
\leq
43C_0(c),
\]
and, for all sufficiently large \(n\), a clipped-shrinkage rule satisfying
\[
\sup_{z\in\mathcal T_K^n}
R_n(\mathcal D^\star,\widehat\tau_n^{\mathrm{sh}};z)
\le
C_0(c)\{n^{-1}-\kappa_c n^{-4/3}\}.
\]
By the definition of \(\rho_n(c)\) in \cref{obj:def:labeled-schedule-game},
\[
\rho_n(c)
\le
C_0(c)\{n^{-1}-\kappa_c n^{-4/3}\},
\]
and hence
\[
d_n(c)\ge C_0(c)\kappa_c n^{-4/3},
\qquad
x_n=s_nd_n(c)\ge C_0(c)\kappa_c ,
\]
for all sufficiently large \(n\), using \(s_n n^{-4/3}=1\). On the other side, \cref{obj:thm:embedded-two-arm-converse} gives, for every \(n\ge8\),
\[
d_n(c)\le 43C_0(c)n^{-4/3},
\qquad
x_n\le 43C_0(c).
\]
Thus, eventually,
\[
x_n\in
\left[C_0(c)\kappa_c,\,43C_0(c)\right],
\qquad
0<x_n.
\]
The displayed liminf--limsup inequalities above give the first assertion of the theorem. % lean: CausalSmith.Experimentation.MultiarmSecondorderMinimaxFrontier.second_order_rate_and_certificate_frontier

\item \textbf{Subsequential limits.}
Put
\[
\mathfrak L_c
=
\left\{y\in\mathbb R:
\text{there is a strictly increasing }\phi:\mathbb N\to\mathbb N
\text{ with }x_{\phi(n)}\to y\right\}.
\]
The eventual containment from the previous step places every tail of \((x_n)\) in the compact interval
\[
I_c=\left[C_0(c)\kappa_c,\,43C_0(c)\right].
\]
Bolzano--Weierstrass applied to a tail subsequence gives \(\mathfrak L_c\neq\varnothing\). Also \(\mathfrak L_c\subseteq I_c\). To see that \(\mathfrak L_c\) is closed, take \(y_j\in\mathfrak L_c\) with \(y_j\to y\). For each \(j\), choose an index \(n_j>n_{j-1}\) such that
\[
|x_{n_j}-y_j|<\frac1j,
\]
possible because \(y_j\) is a subsequential limit. Then \(x_{n_j}\to y\), so \(y\in\mathfrak L_c\). Hence \(\mathfrak L_c\) is a closed subset of the compact interval \(I_c\), and is compact. % lean: CausalSmith.Experimentation.MultiarmSecondorderMinimaxFrontier.subsequentialLimitSet_nonempty_compact

\item \textbf{Regular-variation index and the slowly varying criterion.}
The index assertion is exactly the regular-variation conclusion in \cref{obj:thm:universal-second-order-rate}: if \(a_n>0\) is regularly varying with index \(\beta\), \(C>0\), and
\[
a_n d_n(c)\to C,
\]
then
\[
\beta=\frac43.
\]

It remains to identify when a normalizer of index \(4/3\) exists. Suppose first that \(a_n>0\), \(a_n\) is regularly varying with index \(4/3\), \(C>0\), and \(a_n d_n(c)\to C\). For \(t>0\), set
\[
p_n=a_n d_n(c).
\]
Since \(\lfloor tn\rfloor\to\infty\),
\[
\frac{p_{\lfloor tn\rfloor}}{p_n}\to1.
\]
Moreover,
\[
\frac{s_{\lfloor tn\rfloor}}{s_n}\to t^{4/3},
\qquad
\frac{a_{\lfloor tn\rfloor}}{a_n}\to t^{4/3}.
\]
For all sufficiently large \(n\), the denominators are nonzero, and
\[
\frac{x_{\lfloor tn\rfloor}}{x_n}
=
\frac{s_{\lfloor tn\rfloor}d_{\lfloor tn\rfloor}(c)}
{s_nd_n(c)}
=
\frac{s_{\lfloor tn\rfloor}}{s_n}\,
\frac{p_{\lfloor tn\rfloor}}{p_n}\,
\left(\frac{a_{\lfloor tn\rfloor}}{a_n}\right)^{-1}
\longrightarrow
t^{4/3}\cdot1\cdot t^{-4/3}=1.
\]
Together with the eventual positivity established above, this is the asserted eventual positivity and slow variation of \(x_n\).

Conversely, suppose \(x_n>0\) eventually and
\[
\frac{x_{\lfloor tn\rfloor}}{x_n}\to1
\qquad(t>0).
\]
Choose an integer \(N_+\ge1\) such that \(x_n>0\) for all \(n\ge N_+\), and define a positive sequence \(b_n\) on all of \(\mathbb N\) by assigning arbitrary positive values for \(n<N_+\) and setting
\[
b_n=x_n
\qquad(n\ge N_+).
\]
Since \(\lfloor tn\rfloor\to\infty\), the finite initial assignments are invisible in the ratio, and hence
\[
\frac{b_{\lfloor tn\rfloor}}{b_n}\to1,
\]
so \(b_n\) is regularly varying with index \(0\). Now define a positive sequence \(a_n\) on all of \(\mathbb N\) by assigning arbitrary positive values for \(n<N_+\) and setting
\[
a_n=\frac{s_n}{b_n}
\qquad(n\ge N_+).
\]
This totalization is positive because \(N_+\ge1\), so \(s_n=n^{4/3}>0\) on the displayed tail and \(b_n>0\) everywhere. Since finite initial values do not affect regular variation, and since \(s_n\) is regularly varying with index \(4/3\), the quotient \(a_n=s_n/b_n\) on the tail is regularly varying with index \(4/3\). For all \(n\ge N_+\),
\[
a_n d_n(c)
=
\frac{s_n}{x_n}d_n(c)
=
1.
\]
Thus the desired normalizer exists, with \(C=1\). % lean: CausalSmith.Experimentation.MultiarmSecondorderMinimaxFrontier.second_order_rate_and_certificate_frontier

\item \textbf{Exact rational grid-certificate cluster.}
Let \(q\) be a rational contrast whose real embedding is \(c\). For each \(n\ge1\), apply \cref{obj:thm:rational-contrast-grid-certificate-sandwich} with mesh \(M=n\). This gives rational certificate data
\[
\pi_n,\quad w_n,\quad u_n,\quad \nu_n,\quad \delta_n
\]
such that the exact primal-dual certificate and barycenter identities hold, and such that, writing
\[
\mathsf B_n^q=B_{n,q}(\nu_n),
\qquad
\mathsf U_n^q=U_{n,n}(q;\pi_n,\delta_n),
\qquad
\mathsf T_n^q=u_n=\Lambda_{n,n}(q),
\]
one has
\[
\mathsf B_n^q
\le
\rho_n(c)
\le
\mathsf U_n^q
\le
\mathsf T_n^q
\le
\mathsf B_n^q+\frac{C_0(c)}{4n^2}.
\]
Define, with arbitrary zero-th values,
\[
G_n=s_n\left\{\frac{C_0(c)}{n}-\mathsf T_n^q\right\},
\quad
P_n=s_n\left\{\frac{C_0(c)}{n}-\mathsf U_n^q\right\},
\quad
L_n=s_n\left\{\frac{C_0(c)}{n}-\mathsf B_n^q\right\}.
\]
The same certificate data make \(G_n\) an exact rational grid-certificate cluster for \(q\). The sandwich width satisfies
\[
0\le
s_n(\mathsf T_n^q-\mathsf B_n^q)
\le
s_n\frac{C_0(c)}{4n^2}
=
\frac{C_0(c)}4\,n^{-2/3}
\longrightarrow0.
\]
Since
\[
\mathsf B_n^q\le \rho_n(c)\le \mathsf T_n^q,
\]
we get
\[
|G_n-x_n|
\le
s_n(\mathsf T_n^q-\mathsf B_n^q)\to0.
\]
Similarly, because \(\mathsf U_n^q\le\mathsf T_n^q\) and \(\mathsf B_n^q\le\mathsf U_n^q\),
\[
|P_n-G_n|
=
s_n(\mathsf T_n^q-\mathsf U_n^q)
\le
s_n(\mathsf T_n^q-\mathsf B_n^q)\to0,
\]
and
\[
|L_n-G_n|
=
s_n(\mathsf T_n^q-\mathsf B_n^q)\to0.
\]
These are precisely
\[
G_n-x_n(c)\to0,\qquad
P_n-G_n\to0,\qquad
L_n-G_n\to0.
\] % lean: CausalSmith.Experimentation.MultiarmSecondorderMinimaxFrontier.rational_contrast_grid_certificate_sandwich

\item \textbf{Transferred certificates for real contrasts.}
Apply \cref{obj:thm:real-contrast-grid-certificate-transfer} to the real contrast \(c\). It gives rational contrasts \(c^{(n)}\), certificate values \(B_n,U_n\), and transferred endpoints \(R_n^-,R_n^+\) such that, for every arm \(a\),
\[
c_a^{(n)}\to c_a,
\qquad
n^{5/6}\eta(c,c^{(n)})\to0,
\]
where
\[
\eta(c,c^{(n)})
=
\frac12\sum_{a\in\mathcal A_K}|c_a-c_a^{(n)}|.
\]
For every \(n\ge1\), the transferred certificate supplies an exact rational primal-dual certificate, a grid barycenter estimator, and a rational prior for \(c^{(n)}\), with
\[
B_n=\text{the lower certificate from the prior},
\qquad
U_n=\text{the upper certificate from the grid barycenter},
\]
and
\[
R_n^-=
\left(\max\left\{0,\sqrt{B_n}-\eta(c,c^{(n)})\right\}\right)^2,
\qquad
R_n^+=
\left(\sqrt{U_n}+\eta(c,c^{(n)})\right)^2.
\]
The same transferred certificate gives
\[
R_n^-\le \rho_n(c)\le R_n^+
\]
and, at mesh \(M=n\),
\[
s_n(R_n^+-R_n^-)\to0.
\]
Therefore
\[
0
\le
s_n\left\{\frac{C_0(c)}{n}-R_n^-\right\}-x_n(c)
=
s_n\{\rho_n(c)-R_n^-\}
\le
s_n(R_n^+-R_n^-)\to0,
\]
and
\[
0
\le
x_n(c)-s_n\left\{\frac{C_0(c)}{n}-R_n^+\right\}
=
s_n\{R_n^+-\rho_n(c)\}
\le
s_n(R_n^+-R_n^-)\to0.
\]
Thus
\[
n^{4/3}\left(\frac{C_0(c)}{n}-R_n^-\right)-x_n(c)\to0,
\qquad
n^{4/3}\left(\frac{C_0(c)}{n}-R_n^+\right)-x_n(c)\to0.
\] % lean: CausalSmith.Experimentation.MultiarmSecondorderMinimaxFrontier.real_contrast_grid_certificate_transfer

\item \textbf{The three-arm scalar-score separation.}
For \(K=3\) and \(c^\dagger=(1,-1/2,-1/2)\), \cref{obj:prop:k3-scalar-score-not-minimax-preserving} supplies a full observed arm-label-and-outcome estimator \(\widehat\tau\) under the product \(q^\star\) design \(\mathcal D^\star\) whose worst-case risk over \(\mathcal T_3^3\) is
\[
\frac{511653}{4000000}.
\]
The same result identifies the three-observation scalar sign-score minimax value as
\[
V_3^Z=1-\frac{\sqrt3}{2}
\]
and gives the strict numerical chain
\[
\frac{511653}{4000000}
<
\frac{18213}{136000}
<
1-\frac{\sqrt3}{2}.
\]
Consequently
\[
\sup_{z\in\mathcal T_3^3}
R_3(\mathcal D^\star,\widehat\tau;z)
<
V_3^Z,
\]
as claimed. % lean: CausalSmith.Experimentation.MultiarmSecondorderMinimaxFrontier.k3_scalar_score_not_minimax_preserving
\end{enumerate}
\end{proof}

\begin{proof}[Proof of \cref{obj:thm:attainment-and-k3-certified-converse}]
\begin{enumerate}
\item Suppose first that \(n\ge 1\) and \(|S_c|=2\). Choose the two active arms \(a_+(c)\) and \(a_-(c)\) so that
\[
c_{a_+(c)}=\alpha_c,\qquad c_{a_-(c)}=-\alpha_c,\qquad
\alpha_c=\frac{L_c}{2},
\]
and hence
\[
\alpha_c^2=C_0(c).
\]
For \(z_K=(t_i)_{i=1}^n\in\mathcal T_K^n\), define the active two-arm schedule
\[
\operatorname{act}_c(z_K)_i
=
\bigl(t_{i,a_+(c)},t_{i,a_-(c)}\bigr)\in\mathcal T_2.
\]
For \(z_2=(s_i)_{i=1}^n\in\mathcal T_2^n\), define the sign-schedule embedding \(E_c(z_2)\in\mathcal T_K^n\) by
\[
E_c(z_2)_{i,a_+(c)}=s_{i,1},\qquad
E_c(z_2)_{i,a_-(c)}=s_{i,2},\qquad
E_c(z_2)_{i,a}=0\quad(a\notin S_c).
\]
Then
\[
\tau_c(z_K)
=
\alpha_c\,\tau_{(1,-1)}\{\operatorname{act}_c(z_K)\},
\qquad
\tau_c\{E_c(z_2)\}
=
\alpha_c\,\tau_{(1,-1)}(z_2).
\]
By \cref{obj:thm:embedded-two-arm-converse},
\[
\rho_n(c)=C_0(c)\rho_{n,2}.
\]
% lean: embedded_two_arm_converse

\item Apply \cref{obj:thm:exact-response-type-game} with \(K=2\) and contrast \((1,-1)\). It gives a two-arm orbit procedure \(q_2\) and a prior \(\nu_2\) on \(\mathcal M_{n,2}\) such that
\[
R^{\mathrm{orb}}_{(1,-1)}(q_2;m)\le G_n((1,-1))
\quad(m\in\mathcal M_{n,2}),
\]
and, for every two-arm orbit procedure \(q_2'\),
\[
G_n((1,-1))
\le
\mathbb E_{\nu_2}\!\left[
R^{\mathrm{orb}}_{(1,-1)}(q_2';m)
\right].
\]
The same result identifies \(G_n((1,-1))=\rho_{n,2}\). Let \(p_2\) be the invariant labeled procedure induced by \(q_2\). For \(m\in\mathcal M_{n,2}\), set
\[
\mathrm{Orb}_2(m)
=
\{z_2\in\mathcal T_2^n:\operatorname{counts}(z_2)=m\},
\]
and define the labeled prior
\[
\mathrm{prior}_2(z_2)
=
\frac{\nu_2(\operatorname{counts}(z_2))}
{|\mathrm{Orb}_2(\operatorname{counts}(z_2))|}.
\]
Since \(z_2\in \mathrm{Orb}_2(\operatorname{counts}(z_2))\), the denominator is positive for every schedule in the displayed formula. The invariant correspondence gives
\[
R_n^{(1,-1)}(p_2;z_2)
=
R^{\mathrm{orb}}_{(1,-1)}
\bigl(q_2;\operatorname{counts}(z_2)\bigr),
\]
so
\[
R_n^{(1,-1)}(p_2;z_2)\le \rho_{n,2}
\quad(z_2\in\mathcal T_2^n).
\]
For any two-arm labeled procedure \(p_2'\), symmetrizing \(p_2'\) over unit permutations and passing through the invariant-orbit correspondence yields an orbit procedure whose orbit risk, averaged under \(\nu_2\), is bounded above by the \(\mathrm{prior}_2\)-Bayes risk of \(p_2'\). Therefore
\[
\rho_{n,2}
\le
\mathbb E_{\mathrm{prior}_2}\!\left[
R_n^{(1,-1)}(p_2';z_2)
\right].
\]
% lean: exact_response_type_game

\item Let \(p_K\) be the sign-pair lift of \(p_2\), and let
\[
\mathrm{prior}_K=(E_c)_\#\mathrm{prior}_2,
\qquad
\mathbb E_{\mathrm{prior}_K}F
=
\mathbb E_{\mathrm{prior}_2}\!\left[F\{E_c(z_2)\}\right].
\]
The lifted experiment observes exactly the two active coordinates and multiplies the two-arm estimate by \(\alpha_c\). Hence, for every \(z_K\in\mathcal T_K^n\),
\[
R_n^c(p_K;z_K)
=
\alpha_c^2
R_n^{(1,-1)}
\bigl(p_2;\operatorname{act}_c(z_K)\bigr)
\le
C_0(c)\rho_{n,2}
=
\rho_n(c).
\]
Now fix an arbitrary \(K\)-arm labeled procedure \(p_K'=(\mathcal D_K',\widehat\tau_K')\). Define the induced two-arm procedure \(p_{2\leftarrow K}'\) by the following Rao--Blackwellization. Let \(\widetilde{\mathcal D}_K'\) be the product experiment in which \(\mathfrak a\in\mathcal A_K^n\) is drawn from \(\mathcal D_K'\) and, independently, \(\xi_i\in\{0,1\}\) is fair for each unit. Define the sign-group assignment map \(\phi_c:\mathcal A_K^n\times\{0,1\}^n\to\mathcal A_2^n\) by
\[
\phi_c(\mathfrak a,\xi)_i
=
\begin{cases}
1, & c_{\mathfrak a_i}>0,\\
2, & c_{\mathfrak a_i}<0,\\
1, & c_{\mathfrak a_i}=0\text{ and }\xi_i=1,\\
2, & c_{\mathfrak a_i}=0\text{ and }\xi_i=0.
\end{cases}
\]
Thus the fair bits are used exactly on zero-coefficient assignments. Define the inactive-zero pseudo-observation by
\[
\widetilde y_c(\mathfrak a,y)_i
=
\begin{cases}
y_i, & c_{\mathfrak a_i}\ne 0,\\
0, & c_{\mathfrak a_i}=0,
\end{cases}
\]
where \(0\) is the binary response value assigned to inactive coordinates. For a two-arm observed vector \(y\), set
\[
\zeta_c'(\mathfrak a,\xi;y)
=
\alpha_c^{-1}\widehat\tau_K'\{\mathfrak a,\widetilde y_c(\mathfrak a,y)\}.
\]
Here \(\alpha_c=L_c/2\) and \(\alpha_c>0\), while the \(K\)-arm estimator is clipped to \([-L_c/2,L_c/2]=[-\alpha_c,\alpha_c]\). Hence
\[
\zeta_c'(\mathfrak a,\xi;y)\in[-1,1].
\]
The induced two-arm procedure \(p_{2\leftarrow K}'\) has assignment law \((\phi_c)_\#\widetilde{\mathcal D}_K'\) and estimator
\[
\widehat\tau_{2\leftarrow K}'(\mathfrak b,y)
=
\begin{cases}
\mathbb E_{\widetilde{\mathcal D}_K'}\!\left[
\zeta_c'(\mathfrak a,\xi;y)
\,\middle|\,\phi_c(\mathfrak a,\xi)=\mathfrak b
\right],
&\mathbb P_{\widetilde{\mathcal D}_K'}\{\phi_c=\mathfrak b\}>0,\\[0.5ex]
0,&\mathbb P_{\widetilde{\mathcal D}_K'}\{\phi_c=\mathfrak b\}=0.
\end{cases}
\]
The displayed interval containment makes this a legitimate two-arm procedure. For a fixed \(z_2\in\mathcal T_2^n\), write \(Y_{z_2}(\mathfrak b)\) for the observed outcome vector under two-arm assignment \(\mathfrak b\). Conditional Jensen, applied on each fiber of \(\phi_c\), gives
\[
\begin{aligned}
R_n^{(1,-1)}(p_{2\leftarrow K}';z_2)
&\le
\mathbb E_{\widetilde{\mathcal D}_K'}\!\left[
\left\{
\zeta_c'\bigl(\mathfrak a,\xi;Y_{z_2}(\phi_c(\mathfrak a,\xi))\bigr)
-\tau_{(1,-1)}(z_2)
\right\}^2
\right].
\end{aligned}
\]
For every \((\mathfrak a,\xi)\), the sign-schedule embedding and inactive-zero convention give
\[
\widetilde y_c\{\mathfrak a,Y_{z_2}(\phi_c(\mathfrak a,\xi))\}
=
Y_{E_c(z_2)}(\mathfrak a).
\]
Indeed, a positive-coefficient assigned arm reads the first two-arm coordinate, a negative-coefficient assigned arm reads the second, and an inactive assigned arm has binary response value \(0\) on both sides: \(E_c(z_2)\) fixes inactive coordinates at \(0\), while \(\widetilde y_c\) replaces inactive observations by \(0\). Also,
\[
\begin{aligned}
\tau_c\{E_c(z_2)\}
&=\frac1n\sum_{i=1}^n
\left(
\sum_{a:c_a>0}c_a(z_2)_{i,1}
+
\sum_{a:c_a<0}c_a(z_2)_{i,2}
\right)\\
&=\frac1n\sum_{i=1}^n
\left(\alpha_c(z_2)_{i,1}-\alpha_c(z_2)_{i,2}\right)
=\alpha_c\tau_{(1,-1)}(z_2),
\end{aligned}
\]
where the middle equality uses \(\sum_{a:c_a>0}c_a=\alpha_c\) and \(\sum_{a:c_a<0}c_a=-\alpha_c\). Substituting these two identities into the preceding Jensen bound yields
\[
\begin{aligned}
R_n^{(1,-1)}(p_{2\leftarrow K}';z_2)
&\le
\mathbb E_{\widetilde{\mathcal D}_K'}\!\left[
\alpha_c^{-2}
\left\{
\widehat\tau_K'\bigl(\mathfrak a,Y_{E_c(z_2)}(\mathfrak a)\bigr)
-\tau_c\{E_c(z_2)\}
\right\}^2
\right]\\
&=
\alpha_c^{-2}R_n^c\bigl(p_K';E_c(z_2)\bigr),
\end{aligned}
\]
since \(\widetilde{\mathcal D}_K'\) has \(\mathcal D_K'\) as its \(\mathfrak a\)-marginal.
Integrating this inequality under \(\mathrm{prior}_2\) and using the two-arm Bayes lower bound from the preceding step gives
\[
\rho_{n,2}
\le
\mathbb E_{\mathrm{prior}_2}\!\left[
R_n^{(1,-1)}(p_{2\leftarrow K}';z_2)
\right]
\le
\alpha_c^{-2}
\mathbb E_{\mathrm{prior}_K}\!\left[
R_n^c(p_K';z_K)
\right].
\]
Multiplying by \(\alpha_c^2=C_0(c)\) and using \(\rho_n(c)=C_0(c)\rho_{n,2}\) yields
\[
\rho_n(c)
\le
\mathbb E_{\mathrm{prior}_K}\!\left[
R_n^c(p_K';z_K)
\right].
\]
Together with the definitions of \(p_2,\mathrm{prior}_2,p_K,\mathrm{prior}_K\), this proves the two-active-arm value and saddle-transfer assertion.
% lean: attainment_and_k3_certified_converse

\item For arbitrary active support, \cref{obj:thm:universal-second-order-rate} supplies
\[
C_0(c)\kappa_c
\le
\liminf_{n\to\infty} n^{4/3}d_n(c)
\le
\limsup_{n\to\infty} n^{4/3}d_n(c)
\le
43C_0(c).
\]
It also supplies an integer \(N\) such that, for every \(n\ge N\), the clipped-shrinkage procedure satisfies the contrast-\(c\) risk bound
\[
\sup_{z\in\mathcal T_K^n}
R_n^c(\mathcal D^\star,\widehat\tau_n^{\mathrm{sh}};z)
\le
C_0(c)\left\{n^{-1}-\kappa_c n^{-4/3}\right\}.
\]
These are exactly the second-order scale and shrinkage-attainment assertions.
% lean: universal_second_order_rate

\item By \cref{obj:thm:real-contrast-grid-certificate-transfer}, for every rational nonzero zero-sum contrast \(q\) and every \(n,M\ge 1\), there exist real numbers \(B,U,R^-,R^+\) forming a transferred certificate for \((n,M,c,q)\). The same result supplies rational contrasts \(c^{(n)}\) and real sequences \(B_n,U_n,R_n^-,R_n^+\) such that
\[
c^{(n)}_a\to c_a\quad(a\in\mathcal A_K),
\qquad
n^{5/6}\eta(c,c^{(n)})\to 0,
\]
for every \(n\ge 1\), the tuple \((B_n,U_n,R_n^-,R_n^+)\) is a transferred certificate for \((n,n,c,c^{(n)})\), and
\[
n^{4/3}(R_n^+-R_n^-)\to 0.
\]
Finally, if \(q\) represents \(c\) coefficientwise, the exact rational part of the same result gives, for every \(n\ge 1\), a transferred certificate for \((n,n,c,q)\) satisfying
\[
U-B\le \frac{C_0(c)}{4n^2}.
\]
This proves all real-contrast transfer certificate assertions.
% lean: real_contrast_grid_certificate_transfer

\item Apply \cref{obj:thm:k3-grid-certificate-sandwich} with the given \(n,M\ge 1\), and write the objects it supplies as \(\pi,w,u,\nu,\delta\). For these, \(B_n(\nu)\) is the lower certificate formed from the response-count prior \(\nu\), and \(U_{n,M}\) is the upper certificate formed from \(\pi\) and \(\delta\). The cited result gives the exact primal-dual certificate and barycenter conditions for \(c^\dagger\), and its returned inequalities include
\[
B_n(\nu)\le \rho_n^\dagger,
\qquad
\rho_n^\dagger\le U_{n,M},
\qquad
U_{n,M}\le B_n(\nu)+\frac{1}{4M^2}.
\]
These are exactly the displayed sandwich in the three-arm rational certificate assertion.
% lean: k3_grid_certificate_sandwich
\end{enumerate}
\end{proof}
